%-------------------------------------------------------------------
% Document class and package definitions
%-------------------------------------------------------------------
\documentclass[12pt,a4paper,openright,final,twoside]{cseethesis}

%Included for Gather Purpose only:
%input "thesisreferences.bib"

\usepackage{hyperref}

\makeglossaries

\newacronym{nlp}{NLP}{Natural Language Processing}

\begin{document}

\defaultbibliography{thesisreferences}     %% Change this only.
\defaultbibliographystyle{plain}        %% Could be changed if you like 
                                           %% references typeset differently.
%-------------------------------------------------------------------
% Define title, author, etc.
%-------------------------------------------------------------------
\def\thesistitle{An exploration of word meaning using hyperdimensional computing.}

\begin{figure}[h]
 \centering
 \includegraphics[width=0.5\textwidth]{bg_wall.jpg}
\end{figure}

\def\theauthor{Pedro E.\ Alonso}
\def\theaddress{Dept.\ of Computer Science, Electrical and Space Engineering\\
Lule{\aa} University of Technology\\ Lule{\aa}, Sweden}

\def\supervisors{Marcus Liwicki, Foteini Liwicki, Denis Kleyko}
\def\supervisorstring{Supervisors: \supervisors} % Edit here if you have only one supervisor
\def\dedication{To my past-self...}

% Read abstract and preface from separate files.
% Make sure these exist. See example files.
\def\theabstract{This is an abstract, imported from the file \textit{abstract.tex}.

%Please delete the content below and replace with your abstract. The details below are only provided for information purpose.
% These changes were made by Tosin Adewumi and the template provided.

\section*{Abstract section here monograph...}
No more than 500 words.
Keryword: here!
%Enhancements in this template (from the original by Johan E. Carlson) include:
\begin{enumerate}
    \item Update of the SRT full dept. name on the 1st page in the file cseethesis\_example.
    
    \item A change to defaultbibliographystyle\{\textbf{plain}\} in the cseethesis\_example file instead of the old, which gives blank in-text citation.
    
    \item Inclusion in cseethesis\_example of makeglossaries, the example: \\newacronym\{nlp\}\{NLP\}\{Natural Language Processing\} for handling acronyms and the closing printglossary at the end of the file.

    \item Example usage of the acronym package in the paper1 file: \acrfull{nlp} is the full version of this short version \acrshort{nlp} and the long version \acrlong{nlp}. 
    
    \item Addition of \textbf{newcommand\{\textbackslash makepapertobesubmitted\}} in cseethesis file for manuscripts yet to be submitted.
    
    \item Figure included on the front page in the file cseethesis\_example since the old logo alternative is tricky.
\end{enumerate}}
\def\thepreface{The creation of this template has taken several years, and the shape it is in now would not have been possible with out the patient testers out there. To mention just a few who found and reported bugs, and occasionally even provided bug fixes: Gustav Johansson, Sara Sandberg, Yvonne Aitomäki, Fredrik Hägglund, Jesper Martinsson, Patrik Pääjärvi, and Martin Sehlstedt. To all of those I forgot to mention, please accept my apologies.\\[2.5cm]
Luleå, June 2009\\
Johan E.\ Carlson \\ \\
Improvements by Tosin Adewumi, October 9, 2021.}


% Change here if you want to remove the logo printed on the first page

%\def\thelogo{\includegraphics[width=2.5cm]{bg_wall.jpg}} % old EU logo
\def\thelogo{} % no logo

% The definitions above could be put directly in the function call below,
% but is here defined explicitly, for the purpose of clarity.

\startpreamble
  {\thesistitle}
  {\theauthor}
  {\theaddress}
  {\supervisors}
  {\dedication}
  {\theabstract}
  {\thepreface}
  {\thelogo}

%%%%%%%%%%%%%%%%%%%%%%%%%%%%%%%%%%%%%%%%%%%%%%%%%%%%%%%%%%%%%%%%%%%%
%% Begin Part I
%%%%%%%%%%%%%%%%%%%%%%%%%%%%%%%%%%%%%%%%%%%%%%%%%%%%%%%%%%%%%%%%%%%%
\makepartpage{Part I} 

%% Initialize part containing the thesis introduction chapters
\startchapters
\begin{bibunit}
%------------------------ Start chapter 1 --------------------------
% The \makechapter command takes three arguments
%  1) An abbreviated version of the chapter name,
%     to be used as page header
%  2) String to be added to the table of contents
%  3) The chapter name, possibly split in to lines,
%     as in Chapter 2 below.
%
%  The different arguments can have different line breaks.
%
% The actual contents of the chapter is included by removing the
% comment from the \input line below. Make sure the file
% chapter1.tex exists.
%-------------------------------------------------------------------

\def\myquote{``This report, by its very length, defends itself against the risk of being read."\\[.5\baselineskip] Winston Churchill}

\makechapter[\myquote]{Chapter 1: Title}{Thesis Introduction}{Thesis
Introduction\label{ch1}}
\section{Title: To be filled after}
Current version of the \texttt{cseethesis} document class is: \classversion.\\
Last modification: \classdate\\

\noindent As of version 3.0, the template is no longer backwards compatible.

\subsection{Sub1...About the document class}
This document class was originally created in 2002 when I was working on my own PhD thesis. Since then, many people used it, found bugs (and occasionally even corrected them), and suggested improvements.

The style is tailor-made for the typical types of theses that we write at the department, i.e.\ an introductory part followed by a collection of published or submitted research papers. 

The template supports the use of both \LaTeX\ and pdfLaTeX. If you use the command \texttt{$\backslash$includegraphics} to import your figure and you supply the filename with its extension (e.g.\ .eps, .pdf), compilation should be possible with either one. For this to work, both .eps and .pdf versions of all figures must be available.

The template is totally free to use, modify and distribute, as long as reference to the original author is kept and as long as all files remain in the package. Modified versions can only be distributed if it is clearly mentioned in the document class that modifications have been made and by whom.

The whole package comes AS IS. I will correct bugs every now and then, but other than that, don't expect any support whatsoever.


\subsection{Sub2...About this document}
This document, as well as the actual \LaTeX\ code for it, makes up the documentation on how to use the document class.

Only this chapter contains any readable information. Chapters 2 and 3 are only included as examples of some of the features of the template. The text is nonsense, but the corresponding \LaTeX\ code may be of some use. The same goes for the appended papers, which are only there as examples of a few options of the document class.

Read this chapter carefully. If you have comments on what else should be in here in order to simplify the use of the document class, let me know.

\section{Introduction...Chapters}
\subsection{Defining chapters}
Background info and context
General focus
Thesis statement (research Qs?)
Attract attn.
Why this topic
%In order to add flexibility to the template, a new command, called \texttt{$\backslash$makechapter} is provided. This command takes three mandatory and one optional argument, as 
\begin{center}
	\texttt{$\backslash$makechapter[optional quote]\{page header\}\{toc entry\}\{Chapter title\}}
\end{center}

The reason this is solved like this is to allow for shorter page headers if the chapter name is very long. Also, if the actual chapter heading needs to be manually split in several lines (if the automatic splitting does not look so good), the table of contents (toc) entry might have to be defined differently. Note that normally, the last three arguments can be the same. 

The use of an optional quote as an introduction to the chapter is demonstrated in this chapter. It can just as well be left out, which is demonstrated in this document (see the code).

\subsection{Importing chapter contents}
The sub-documents containing the chapters should start directly, i.e.\ they must not contain any $\backslash$begin\{document\} or $\backslash$end\{document\} tags.

See this file, \textit{chapter1.tex} for details.

\section{Methodology section...How to append papers}
data collection analysis and interpretation how? Here!
Description of the process

%Papers are included using the \texttt{$\backslash$input} command, just as with chapters. You have to typeset paper title, authors, and abstract manually. See the example papers accompanying this document for an example.

To make the separator sheet preceding each paper, use one of the following commands:
\begin{itemize}
	\item \texttt{$\backslash$makepaper} -- Published paper.
	\item \texttt{$\backslash$makepaperaccepted} -- Accepted, not yet published paper.
	\item \texttt{$\backslash$makepapersubmitted} -- Submitted, not yet accepted paper.
	\item \texttt{$\backslash$makepapertobesubmitted} -- Not yet submitted paper.
\end{itemize}

See code for this example document for examples on how to use.

\section{Literature Review... Cross-references}
Overview of major or important works for this topic
Cite
Find gap
Fill gap
%All labels throughout the thesis have to be unique. If the same
%equation or figure shows up twice e.g.\ a figure used
%both in the introduction and one of the included papers), it has
%to be given different labels. Referring to labels is done as
usual.

A simple trick to make sure this is the case and that will also help you keep track of all labels you used is to use the following naming convention:
\begin{itemize}
	\item \texttt{ch1:fig:labelname}, \texttt{ch1:tab:labelname}, \texttt{ch1:eq:labelname}, etc.\ all denote figures, tables and equations in Chapter 1.
	\item \texttt{paperA:fig:labelname}, \texttt{paperA:tab:labelname}, \texttt{paperA:eq:labelname}, etc.\ all denote figures, tables and equations in Paper A.
\end{itemize}

For existing text, e.g.\ papers, this is easily achieved by a simple search-and-replace operation on the string \texttt{$\backslash$label\{}. Any text editor will do that for you!

\section{Results Section.... Appendices\label{sec:app}}
What did I find?
What did I not find?
What I did find that was not expected?
%It is possible to have any number of appendices for each chapter.
%Simply type \texttt{$\backslash$appendix} before the first appendix,
%which is then a normal \texttt{section}, but numbered differently.

To add appendices to papers, use the
\texttt{$\backslash$paperappendix} command

\section{Dicussion Section... Including bibliography lists\label{sec:bib}}
Interpret results
Data supports goals?
Contribution or new?
Limitation (scope)
%In a thesis one might want several separate bibliographies. For
%example, one for the first part, and then separate bibliographies
%for each of the included papers.

This is solved using the \texttt{bibunits} package together with a
slight work-around in this template. For the first part of the
thesis, there is only one bibliography list, typeset like a chapter
(see this example document). In the papers, the bibliography lists
are typeset as un-numbered sections. See this file
\texttt{cseethesis\_example.tex} and \texttt{paper1.tex} for examples
how to place the bibliographies. 

Note that the command
\begin{itemize}
	\item \texttt{$\backslash$makebib} is used in Part I, to typeset the reference list in the thesis introduction.
	\item \texttt{$\backslash$putbib} is used in the papers in Part II.
\end{itemize}

\section{Conclusion Section... How to compile your project}
Finally, you probably like to know how to build your project to a
final PDF or PostScript file. Start by verifying that you can
compile this document. This is how it goes:
%
\begin{enumerate}
    \item Run \LaTeX\ (or pdfLaTex) once.
    \item Then run BibTeX on all the
    bu<i> files.
    \item Run \LaTeX\ twice more, to build the final DVI document
    (or pdfLaTeX if you want a PDF file).
\end{enumerate}

The above steps are easily collected in a script or batch file. See
the files \texttt{make.bat} and \texttt{compilebibunits.bat} for
examples.

\section{If I'm here is completed... Revision history}
By following :  \href{https://www.slideshare.net/NasreddineELGUEZAR/how-to-write-a-research-monograph-basics}{presentation} The template has evolved during several years and the exact
revisions are not clear to anyone. Starting from version 1.6,
however, the changes are more well-documented. This example
document will always support only the latest release of the
template. Below is a list of the revisions made to the document
class (and when applicable, the example document):%
\begin{itemize}
	 \item Version 3.1, September 1, 2010
		\begin{itemize}
			\item Fixed bug related to appendix numbering.
			\item Fixed page numbering of ``Part" pages.
			\item Added a comment at top of cseethesis\_example.tex, for improved compatibility with some editors.
		\end{itemize}
	 \item Version 3.0, June 7, 2009
		\begin{itemize}
			\item Fixed bug related to page headers in chapters containing no subsections.
		   \item Removed the EU class option and replaced with a logo argument to the preamble. See the code of this document for an example.
			\item The template is \textbf{no longer compatible with previous versions}.
         \item Removed the definition of boldface Greek letters from the document class, since this is not the proper place for that.
		\end{itemize}
  	 \item Version 2.5, March 5 2009
		\begin{itemize}
         \item Renamed the template \textit{cseethesis}. It is a continuation of the project initially called \textit{eisthesis}, but since its use has spread I decided to change the name.
			\item Various minor bug fixes.
			\item Update of example document, making examples of additions and revisions in recent versions of the template.
		\end{itemize}
	 \item Version 2.36: Added "Part" to the table of contents.
    \item Version 2.35:
        \begin{itemize}
				\item Fixed a bug regarding the page headers in the "appended papers part".
				\item Fixed a bug causing the section numbering to be wrong in a chapter succeeding a chapter containing appendices.
			   \item Added a chapter 3 in this example document, illustrating how to handle page headers for chapters without any sections. See the code at the top of \texttt{chapter3.tex}.
		  \end{itemize}
    \item Version 2.3:
        \begin{itemize}
            \item Added the commands \texttt{$\backslash$appendix}
            and \texttt{$\backslash$paperappendix}, see section
            \ref{sec:app}.
            \item The bibliography list is now typeset similar to a
            chapter in Part I, and as an un-numbered section in Part
            II. This required the use of separate commands for
            including the lists (see Sec.~\ref{sec:bib})
            \item Typesetting fixes for the table of contents page.
            As a consequence, the package \texttt{titletoc} is now
            required.
            \item Minor other code cleanup and bug fixes.
        \end{itemize}
    \item Version 2.25:
    \begin{itemize}
        \item Fixed a minor bug that used to generate a warning
        message regarding font shapes in the page headers.
    \end{itemize}
    \item Version 2.2:
    \begin{itemize}
        \item pdf and eps class options removed. The document class
        compiles with either pdfLaTeX or \LaTeX.
        \item The reference lists in papers are now typeset in the
        same way as in Part I of the thesis.
        \item Some minor adjustments of page header heights.
    \end{itemize}
    \item Version 2.1:
    \begin{itemize}
        \item Cross-references to chapters now work like they
        should. See the main document of this example.
        \item Major bug fixes to BibTeX reference lists, table of
        contents generation etc.
        \item The only change the user has to do is to use the new \texttt{$\backslash$makebib} instead of
        bibunits' \texttt{$\backslash$putbib}.
    \end{itemize}
    \item Version 2.0:
    \begin{itemize}
        \item Support for EU logotype on main page. The files
        \texttt{eu1\_f\_eng.pdf} and\\ \texttt{eu1\_f\_eng.eps} must be
        placed in the same directory as the document class.
        \item Support for pdf class options.
        \item Update of the \texttt{$\backslash$makechapter}
        command. It now requires three arguments. See main
        document for example.
        \item Support for BibTeX, using the \texttt{bibunits.sty}
        package.
    \end{itemize}
    \item Version 1.6: Various bug fixes to figure spacing etc.
\end{itemize}


\makechapter{Chapter 2: Methodology Chapter}{Table of contents entry}{Methods used\\ \label{ch2}}
\section{First section of the second chapter}
This is the text of the second chapter. \cite{*}

\appendix
\section{This is an appendix section}
Text of the appendix

\subsection{Subsection 1}
Yet some text, and an equation
\begin{equation}
    \textnormal{abs}\left(e^{j\pi}\right)=\,?
\end{equation}

\subsection{Subsection 2}
And then some...

\section{This is another appendix section}
This section concludes the appendix.


\makechapter{Literature Review}{Table of contents}{Literature Surveyed\\ \label{ch3}}
This is a chapter with no sections, only here in order to test the document template. Please ignore the rest of this.
\vspace*{24pt}
dkhgfdsh fdsjhgiepy kdslkfds ewiuyfe fkjfdsohew dskjhfd fduew dsk di e sd lkjd dslkfd oiew sao cdkåwq acslkgå sdspo dsjdspe dkfo dkhgfdsh fdsjhgiepy kdslkfds ewiuyfe fkjfdsohew dskjhfd fduew dsk di e sd lkjd dslkfd oiew sao cdkåwq acslkgå sdspo dsjdspe dkfo dkhgfdsh fdsjhgiepy kdslkfds ewiuyfe fkjfdsohew dskjhfd fduew dsk di e sd lkjd dslkfd oiew sao cdkåwq acslkgå sdspo dsjdspe dkfo

dkhgfdsh fdsjhgiepy kdslkfds ewiuyfe fkjfdsohew dskjhfd fduew dsk di e sd lkjd dslkfd oiew sao cdkåwq acslkgå sdspo dsjdspe dkfodkhgfdsh fdsjhgiepy kdslkfds ewiuyfe fkjfdsohew dskjhfd fduew dsk di e sd lkjd dslkfd oiew sao cdkåwq acslkgå sdspo dsjdspe dkfo dkhgfdsh fdsjhgiepy kdslkfds ewiuyfe fkjfdsohew dskjhfd fduew dsk di e sd lkjd dslkfd oiew sao cdkåwq acslkgå sdspo dsjdspe dkfo dkhgfdsh fdsjhgiepy kdslkfds ewiuyfe fkjfdsohew dskjhfd fduew dsk di e sd lkjd dslkfd oiew sao cdkåwq acslkgå sdspo dsjdspe dkfo

dkhgfdsh fdsjhgiepy kdslkfds ewiuyfe fkjfdsohew dskjhfd fduew dsk di e sd lkjd dslkfd oiew sao cdkåwq acslkgå sdspo dsjdspe dkfodkhgfdsh fdsjhgiepy kdslkfds ewiuyfe fkjfdsohew dskjhfd fduew dsk di e sd lkjd dslkfd oiew sao cdkåwq acslkgå sdspo dsjdspe dkfo dkhgfdsh fdsjhgiepy kdslkfds ewiuyfe fkjfdsohew dskjhfd fduew dsk di e sd lkjd dslkfd oiew sao cdkåwq acslkgå sdspo dsjdspe dkfo dkhgfdsh fdsjhgiepy kdslkfds ewiuyfe fkjfdsohew dskjhfd fduew dsk di e sd lkjd dslkfd oiew sao cdkåwq acslkgå sdspo dsjdspe dkfo dkhgfdsh fdsjhgiepy kdslkfds ewiuyfe fkjfdsohew dskjhfd fduew dsk di e sd lkjd  dslkfd oiew sao cdkåwq acslkgå sdspo dsjdspe dkfodkhgfdsh fdsjhgiepy kdslkfds ewiuyfe fkjfdsohew dskjhfd fduew dsk di e sd lkjd dslkfd oiew sao cdkåwq acslkgå sdspo dsjdspe dkfo
dkhgfdsh fdsjhgiepy kdslkfds ewiuyfe fkjfdsohew dskjhfd fduew dsk di e sd lkjd dslkfd oiew sao cdkåwq acslkgå sdspo dsjdspe dkfo
dkhgfdsh fdsjhgiepy kdslkfds ewiuyfe fkjfdsohew dskjhfd fduew dsk di e sd lkjd dslkfd oiew sao cdkåwq acslkgå sdspo dsjdspe dkfo
dkhgfdsh fdsjhgiepy kdslkfds ewiuyfe fkjfdsohew dskjhfd fduew dsk di e sd lkjd dslkfd oiew sao cdkåwq acslkgå sdspo dsjdspe dkfo
dkhgfdsh fdsjhgiepy kdslkfds ewiuyfe fkjfdsohew dskjhfd fduew dsk di e sd lkjd dslkfd oiew sao cdkåwq acslkgå sdspo dsjdspe dkfo
dkhgfdsh fdsjhgiepy kdslkfds ewiuyfe fkjfdsohew dskjhfd fduew dsk di e sd lkjd dslkfd oiew sao cdkåwq acslkgå sdspo dsjdspe dkfo

dkhgfdsh fdsjhgiepy kdslkfds ewiuyfe fkjfdsohew dskjhfd fduew dsk di e sd lkjd dslkfd oiew sao cdkåwq acslkgå sdspo dsjdspe dkfo
dkhgfdsh fdsjhgiepy kdslkfds ewiuyfe fkjfdsohew dskjhfd fduew dsk di e sd lkjd dslkfd oiew sao cdkåwq acslkgå sdspo dsjdspe dkfo

dkhgfdsh fdsjhgiepy kdslkfds ewiuyfe fkjfdsohew dskjhfd fduew dsk di e sd lkjd dslkfd oiew sao cdkåwq acslkgå sdspo dsjdspe dkfo
 
dkhgfdsh fdsjhgiepy kdslkfds ewiuyfe fkjfdsohew dskjhfd fduew dsk di e sd lkjd dslkfd oiew sao cdkåwq acslkgå sdspo dsjdspe dkfodkhgfdsh fdsjhgiepy kdslkfds ewiuyfe fkjfdsohew dskjhfd fduew dsk di e sd lkjd dslkfd oiew sao cdkåwq acslkgå sdspo dsjdspe dkfodkhgfdsh fdsjhgiepy kdslkfds ewiuyfe fkjfdsohew dskjhfd fduew dsk di e sd lkjd dslkfd oiew sao cdkåwq acslkgå sdspo dsjdspe dkfodkhgfdsh fdsjhgiepy kdslkfds ewiuyfe fkjfdsohew dskjhfd fduew dsk di e sd lkjd dslkfd oiew sao cdkåwq acslkgå sdspo dsjdspe dkfodkhgfdsh fdsjhgiepy kdslkfds ewiuyfe fkjfdsohew dskjhfd fduew dsk di e sd lkjd dslkfd oiew sao cdkåwq acslkgå sdspo dsjdspe dkfodkhgfdsh fdsjhgiepy kdslkfds ewiuyfe fkjfdsohew dskjhfd fduew dsk di e sd lkjd dslkfd oiew sao cdkåwq acslkgå sdspo dsjdspe dkfodkhgfdsh fdsjhgiepy kdslkfds ewiuyfe fkjfdsohew dskjhfd fduew dsk di e sd lkjd dslkfd oiew sao cdkåwq acslkgå sdspo dsjdspe dkfodkhgfdsh fdsjhgiepy kdslkfds ewiuyfe fkjfdsohew dskjhfd fduew dsk di e sd lkjd dslkfd oiew sao cdkåwq acslkgå sdspo dsjdspe dkfodkhgfdsh fdsjhgiepy kdslkfds ewiuyfe fkjfdsohew dskjhfd fduew dsk di e sd lkjd dslkfd oiew sao cdkåwq acslkgå sdspo dsjdspe dkfodkhgfdsh fdsjhgiepy kdslkfds ewiuyfe fkjfdsohew dskjhfd fduew dsk di e sd lkjd dslkfd oiew sao cdkåwq acslkgå sdspo dsjdspe dkfodkhgfdsh fdsjhgiepy kdslkfds ewiuyfe fkjfdsohew dskjhfd fduew dsk di e sd lkjd dslkfd oiew sao cdkåwq acslkgå sdspo dsjdspe dkfo
dkhgfdsh fdsjhgiepy kdslkfds ewiuyfe fkjfdsohew dskjhfd fduew dsk di e sd lkjd dslkfd oiew sao cdkåwq acslkgå sdspo dsjdspe dkfodkhgfdsh fdsjhgiepy kdslkfds ewiuyfe fkjfdsohew dskjhfd fduew dsk di e sd lkjd dslkfd oiew sao cdkåwq acslkgå sdspo dsjdspe dkfo

dkhgfdsh fdsjhgiepy kdslkfds ewiuyfe fkjfdsohew dskjhfd fduew dsk di e sd lkjd dslkfd oiew sao cdkåwq acslkgå sdspo dsjdspe dkfodkhgfdsh fdsjhgiepy kdslkfds ewiuyfe fkjfdsohew dskjhfd fduew dsk di e sd lkjd dslkfd oiew sao cdkåwq acslkgå sdspo dsjdspe dkfodkhgfdsh fdsjhgiepy kdslkfds ewiuyfe fkjfdsohew dskjhfd fduew dsk di e sd lkjd dslkfd oiew sao cdkåwq acslkgå sdspo dsjdspe dkfodkhgfdsh fdsjhgiepy kdslkfds ewiuyfe fkjfdsohew dskjhfd fduew dsk di e sd lkjd dslkfd oiew sao cdkåwq acslkgå sdspo dsjdspe dkfo
dkhgfdsh fdsjhgiepy kdslkfds ewiuyfe fkjfdsohew dskjhfd fduew dsk di e sd lkjd dslkfd oiew sao cdkåwq acslkgå sdspo dsjdspe dkfo

dkhgfdsh fdsjhgiepy kdslkfds ewiuyfe fkjfdsohew dskjhfd fduew dsk di e sd lkjd dslkfd oiew sao cdkåwq acslkgå sdspo dsjdspe dkfodkhgfdsh fdsjhgiepy kdslkfds ewiuyfe fkjfdsohew dskjhfd fduew dsk di e sd lkjd dslkfd oiew sao cdkåwq acslkgå sdspo dsjdspe dkfodkhgfdsh fdsjhgiepy kdslkfds ewiuyfe fkjfdsohew dskjhfd fduew dsk di e sd lkjd dslkfd oiew sao cdkåwq acslkgå sdspo dsjdspe dkfodkhgfdsh fdsjhgiepy kdslkfds ewiuyfe fkjfdsohew dskjhfd fduew dsk di e sd lkjd dslkfd oiew sao cdkåwq acslkgå sdspo dsjdspe dkfodkhgfdsh fdsjhgiepy kdslkfds ewiuyfe fkjfdsohew dskjhfd fduew dsk di e sd lkjd dslkfd oiew sao cdkåwq acslkgå sdspo dsjdspe dkfo

dkhgfdsh fdsjhgiepy kdslkfds ewiuyfe fkjfdsohew dskjhfd fduew dsk di e sd lkjd dslkfd oiew sao cdkåwq acslkgå sdspo dsjdspe dkfodkhgfdsh fdsjhgiepy kdslkfds ewiuyfe fkjfdsohew dskjhfd fduew dsk di e sd lkjd dslkfd oiew sao cdkåwq acslkgå sdspo dsjdspe dkfodkhgfdsh fdsjhgiepy kdslkfds ewiuyfe fkjfdsohew dskjhfd fduew dsk di e sd lkjd dslkfd oiew sao cdkåwq acslkgå sdspo dsjdspe dkfodkhgfdsh fdsjhgiepy kdslkfds ewiuyfe fkjfdsohew dskjhfd fduew dsk di e sd lkjd dslkfd oiew sao cdkåwq acslkgå sdspo dsjdspe dkfodkhgfdsh fdsjhgiepy kdslkfds ewiuyfe fkjfdsohew dskjhfd fduew dsk di e sd lkjd dslkfd oiew sao cdkåwq acslkgå sdspo dsjdspe dkfodkhgfdsh fdsjhgiepy kdslkfds ewiuyfe fkjfdsohew dskjhfd fduew dsk di e sd lkjd dslkfd oiew sao cdkåwq acslkgå sdspo dsjdspe dkfodkhgfdsh fdsjhgiepy kdslkfds ewiuyfe fkjfdsohew dskjhfd fduew dsk di e sd lkjd dslkfd oiew sao cdkåwq acslkgå sdspo dsjdspe dkfo
dkhgfdsh fdsjhgiepy kdslkfds ewiuyfe fkjfdsohew dskjhfd fduew dsk di e sd lkjd dslkfd oiew sao cdkåwq acslkgå sdspo dsjdspe dkfodkhgfdsh fdsjhgiepy kdslkfds ewiuyfe fkjfdsohew dskjhfd fduew dsk di e sd lkjd dslkfd oiew sao cdkåwq acslkgå sdspo dsjdspe dkfo

dkhgfdsh fdsjhgiepy kdslkfds ewiuyfe fkjfdsohew dskjhfd fduew dsk di e sd lkjd dslkfd oiew sao cdkåwq acslkgå sdspo dsjdspe dkfodkhgfdsh fdsjhgiepy kdslkfds ewiuyfe fkjfdsohew dskjhfd fduew dsk di e sd lkjd dslkfd oiew sao cdkåwq acslkgå sdspo dsjdspe dkfodkhgfdsh fdsjhgiepy kdslkfds ewiuyfe fkjfdsohew dskjhfd fduew dsk di e sd lkjd dslkfd oiew sao cdkåwq acslkgå sdspo dsjdspe dkfodkhgfdsh fdsjhgiepy kdslkfds ewiuyfe fkjfdsohew dskjhfd fduew dsk di e sd lkjd dslkfd oiew sao cdkåwq acslkgå sdspo dsjdspe dkfodkhgfdsh fdsjhgiepy kdslkfds ewiuyfe fkjfdsohew dskjhfd fduew dsk di e sd lkjd dslkfd oiew sao cdkåwq acslkgå sdspo dsjdspe dkfo
dkhgfdsh fdsjhgiepy kdslkfds ewiuyfe fkjfdsohew dskjhfd fduew dsk di e sd lkjd dslkfd oiew sao cdkåwq acslkgå sdspo dsjdspe dkfodkhgfdsh fdsjhgiepy kdslkfds ewiuyfe fkjfdsohew dskjhfd fduew dsk di e sd lkjd dslkfd oiew sao cdkåwq acslkgå sdspo dsjdspe dkfo

dkhgfdsh fdsjhgiepy kdslkfds ewiuyfe fkjfdsohew dskjhfd fduew dsk di e sd lkjd dslkfd oiew sao cdkåwq acslkgå sdspo dsjdspe dkfodkhgfdsh fdsjhgiepy kdslkfds ewiuyfe fkjfdsohew dskjhfd fduew dsk di e sd lkjd dslkfd oiew sao cdkåwq acslkgå sdspo dsjdspe dkfodkhgfdsh fdsjhgiepy kdslkfds ewiuyfe fkjfdsohew dskjhfd fduew dsk di e sd lkjd dslkfd oiew sao cdkåwq acslkgå sdspo dsjdspe dkfo

dkhgfdsh fdsjhgiepy kdslkfds ewiuyfe fkjfdsohew dskjhfd fduew dsk di e sd lkjd dslkfd oiew sao cdkåwq acslkgå sdspo dsjdspe dkfodkhgfdsh fdsjhgiepy kdslkfds ewiuyfe fkjfdsohew dskjhfd fduew dsk di e sd lkjd dslkfd oiew sao cdkåwq acslkgå sdspo dsjdspe dkfodkhgfdsh fdsjhgiepy kdslkfds ewiuyfe fkjfdsohew dskjhfd fduew dsk di e sd lkjd dslkfd oiew sao cdkåwq acslkgå sdspo dsjdspe dkfodkhgfdsh fdsjhgiepy kdslkfds ewiuyfe fkjfdsohew dskjhfd fduew dsk di e sd lkjd dslkfd oiew sao cdkåwq acslkgå sdspo dsjdspe dkfodkhgfdsh fdsjhgiepy kdslkfds ewiuyfe fkjfdsohew dskjhfd fduew dsk di e sd lkjd dslkfd oiew sao cdkåwq acslkgå sdspo dsjdspe dkfodkhgfdsh fdsjhgiepy kdslkfds ewiuyfe fkjfdsohew dskjhfd fduew dsk di e sd lkjd dslkfd oiew sao cdkåwq acslkgå sdspo dsjdspe dkfodkhgfdsh fdsjhgiepy kdslkfds ewiuyfe fkjfdsohew dskjhfd fduew dsk di e sd lkjd dslkfd oiew sao cdkåwq acslkgå sdspo dsjdspe dkfo
dkhgfdsh fdsjhgiepy kdslkfds ewiuyfe fkjfdsohew dskjhfd fduew dsk di e sd lkjd dslkfd oiew sao cdkåwq acslkgå sdspo dsjdspe dkfodkhgfdsh fdsjhgiepy kdslkfds ewiuyfe fkjfdsohew dskjhfd fduew dsk di e sd lkjd dslkfd oiew sao cdkåwq acslkgå sdspo dsjdspe dkfo

dkhgfdsh fdsjhgiepy kdslkfds ewiuyfe fkjfdsohew dskjhfd fduew dsk di e sd lkjd dslkfd oiew sao cdkåwq acslkgå sdspo dsjdspe dkfodkhgfdsh fdsjhgiepy kdslkfds ewiuyfe fkjfdsohew dskjhfd fduew dsk di e sd lkjd dslkfd oiew sao cdkåwq acslkgå sdspo dsjdspe dkfodkhgfdsh fdsjhgiepy kdslkfds ewiuyfe fkjfdsohew dskjhfd fduew dsk di e sd lkjd dslkfd oiew sao cdkåwq acslkgå sdspo dsjdspe dkfodkhgfdsh fdsjhgiepy kdslkfds ewiuyfe fkjfdsohew dskjhfd fduew dsk di e sd lkjd dslkfd oiew sao cdkåwq acslkgå sdspo dsjdspe dkfo
dkhgfdsh fdsjhgiepy kdslkfds ewiuyfe fkjfdsohew dskjhfd fduew dsk di e sd lkjd dslkfd oiew sao cdkåwq acslkgå sdspo dsjdspe dkfodkhgfdsh fdsjhgiepy kdslkfds ewiuyfe fkjfdsohew dskjhfd fduew dsk di e sd lkjd dslkfd oiew sao cdkåwq acslkgå sdspo dsjdspe dkfodkhgfdsh fdsjhgiepy kdslkfds ewiuyfe fkjfdsohew dskjhfd fduew dsk di e sd lkjd dslkfd oiew sao cdkåwq acslkgå sdspo dsjdspe dkfodkhgfdsh fdsjhgiepy kdslkfds ewiuyfe fkjfdsohew dskjhfd fduew dsk di e sd lkjd dslkfd oiew sao cdkåwq acslkgå sdspo dsjdspe dkfo
dkhgfdsh fdsjhgiepy kdslkfds ewiuyfe fkjfdsohew dskjhfd fduew dsk di e sd lkjd dslkfd oiew sao cdkåwq acslkgå sdspo dsjdspe dkfodkhgfdsh fdsjhgiepy kdslkfds ewiuyfe fkjfdsohew dskjhfd fduew dsk di e sd lkjd dslkfd oiew sao cdkåwq acslkgå sdspo dsjdspe dkfodkhgfdsh fdsjhgiepy kdslkfds ewiuyfe fkjfdsohew dskjhfd fduew dsk di e sd lkjd dslkfd oiew sao cdkåwq acslkgå sdspo dsjdspe dkfodkhgfdsh fdsjhgiepy kdslkfds ewiuyfe fkjfdsohew dskjhfd fduew dsk di e sd lkjd dslkfd oiew sao cdkåwq acslkgå sdspo dsjdspe dkfo
dkhgfdsh fdsjhgiepy kdslkfds ewiuyfe fkjfdsohew dskjhfd fduew dsk di e sd lkjd dslkfd oiew sao cdkåwq acslkgå sdspo dsjdspe dkfodkhgfdsh fdsjhgiepy kdslkfds ewiuyfe fkjfdsohew dskjhfd fduew dsk di e sd lkjd dslkfd oiew sao cdkåwq acslkgå sdspo dsjdspe dkfodkhgfdsh fdsjhgiepy kdslkfds ewiuyfe fkjfdsohew dskjhfd fduew dsk di e sd lkjd dslkfd oiew sao cdkåwq acslkgå sdspo dsjdspe dkfo

dkhgfdsh fdsjhgiepy kdslkfds ewiuyfe fkjfdsohew dskjhfd fduew dsk di e sd lkjd dslkfd oiew sao cdkåwq acslkgå sdspo dsjdspe dkfodkhgfdsh fdsjhgiepy kdslkfds ewiuyfe fkjfdsohew dskjhfd fduew dsk di e sd lkjd dslkfd oiew sao cdkåwq acslkgå sdspo dsjdspe dkfodkhgfdsh fdsjhgiepy kdslkfds ewiuyfe fkjfdsohew dskjhfd fduew dsk di e sd lkjd dslkfd oiew sao cdkåwq acslkgå sdspo dsjdspe dkfodkhgfdsh fdsjhgiepy kdslkfds ewiuyfe fkjfdsohew dskjhfd fduew dsk di e sd lkjd dslkfd oiew sao cdkåwq acslkgå sdspo dsjdspe dkfodkhgfdsh fdsjhgiepy kdslkfds ewiuyfe fkjfdsohew dskjhfd fduew dsk di e sd lkjd dslkfd oiew sao cdkåwq acslkgå sdspo dsjdspe dkfodkhgfdsh fdsjhgiepy kdslkfds ewiuyfe fkjfdsohew dskjhfd fduew dsk di e sd lkjd dslkfd oiew sao cdkåwq acslkgå sdspo dsjdspe dkfodkhgfdsh fdsjhgiepy kdslkfds ewiuyfe fkjfdsohew dskjhfd fduew dsk di e sd lkjd dslkfd oiew sao cdkåwq acslkgå sdspo dsjdspe dkfodkhgfdsh fdsjhgiepy kdslkfds ewiuyfe fkjfdsohew dskjhfd fduew dsk di e sd lkjd dslkfd oiew sao cdkåwq acslkgå sdspo dsjdspe dkfodkhgfdsh fdsjhgiepy kdslkfds ewiuyfe fkjfdsohew dskjhfd fduew dsk di e sd lkjd dslkfd oiew sao cdkåwq acslkgå sdspo dsjdspe dkfodkhgfdsh fdsjhgiepy kdslkfds ewiuyfe fkjfdsohew dskjhfd fduew dsk di e sd lkjd dslkfd oiew sao cdkåwq acslkgå sdspo dsjdspe dkfodkhgfdsh fdsjhgiepy kdslkfds ewiuyfe fkjfdsohew dskjhfd fduew dsk di e sd lkjd dslkfd oiew sao cdkåwq acslkgå sdspo dsjdspe dkfo
dkhgfdsh fdsjhgiepy kdslkfds ewiuyfe fkjfdsohew dskjhfd fduew dsk di e sd lkjd dslkfd oiew sao cdkåwq acslkgå sdspo dsjdspe dkfodkhgfdsh fdsjhgiepy kdslkfds ewiuyfe fkjfdsohew dskjhfd fduew dsk di e sd lkjd dslkfd oiew sao cdkåwq acslkgå sdspo dsjdspe dkfo



\makechapter{Results Chapter}{Table of contents}{Results obtained \\ \label{ch4}}
\section{Title: To be filled after}
Current version of the \texttt{cseethesis} document class is: \classversion.\\
Last modification: \classdate\\

\noindent As of version 3.0, the template is no longer backwards compatible.

\subsection{Sub1...About the document class}
This document class was originally created in 2002 when I was working on my own PhD thesis. Since then, many people used it, found bugs (and occasionally even corrected them), and suggested improvements.

The style is tailor-made for the typical types of theses that we write at the department, i.e.\ an introductory part followed by a collection of published or submitted research papers. 

The template supports the use of both \LaTeX\ and pdfLaTeX. If you use the command \texttt{$\backslash$includegraphics} to import your figure and you supply the filename with its extension (e.g.\ .eps, .pdf), compilation should be possible with either one. For this to work, both .eps and .pdf versions of all figures must be available.

The template is totally free to use, modify and distribute, as long as reference to the original author is kept and as long as all files remain in the package. Modified versions can only be distributed if it is clearly mentioned in the document class that modifications have been made and by whom.

The whole package comes AS IS. I will correct bugs every now and then, but other than that, don't expect any support whatsoever.


\subsection{Sub2...About this document}
This document, as well as the actual \LaTeX\ code for it, makes up the documentation on how to use the document class.

Only this chapter contains any readable information. Chapters 2 and 3 are only included as examples of some of the features of the template. The text is nonsense, but the corresponding \LaTeX\ code may be of some use. The same goes for the appended papers, which are only there as examples of a few options of the document class.

Read this chapter carefully. If you have comments on what else should be in here in order to simplify the use of the document class, let me know.

\section{Introduction...Chapters}
\subsection{Defining chapters}
Background info and context
General focus
Thesis statement (research Qs?)
Attract attn.
Why this topic
%In order to add flexibility to the template, a new command, called \texttt{$\backslash$makechapter} is provided. This command takes three mandatory and one optional argument, as 
\begin{center}
	\texttt{$\backslash$makechapter[optional quote]\{page header\}\{toc entry\}\{Chapter title\}}
\end{center}

The reason this is solved like this is to allow for shorter page headers if the chapter name is very long. Also, if the actual chapter heading needs to be manually split in several lines (if the automatic splitting does not look so good), the table of contents (toc) entry might have to be defined differently. Note that normally, the last three arguments can be the same. 

The use of an optional quote as an introduction to the chapter is demonstrated in this chapter. It can just as well be left out, which is demonstrated in this document (see the code).

\subsection{Importing chapter contents}
The sub-documents containing the chapters should start directly, i.e.\ they must not contain any $\backslash$begin\{document\} or $\backslash$end\{document\} tags.

See this file, \textit{chapter1.tex} for details.

\section{Methodology section...How to append papers}
data collection analysis and interpretation how? Here!
Description of the process

%Papers are included using the \texttt{$\backslash$input} command, just as with chapters. You have to typeset paper title, authors, and abstract manually. See the example papers accompanying this document for an example.

To make the separator sheet preceding each paper, use one of the following commands:
\begin{itemize}
	\item \texttt{$\backslash$makepaper} -- Published paper.
	\item \texttt{$\backslash$makepaperaccepted} -- Accepted, not yet published paper.
	\item \texttt{$\backslash$makepapersubmitted} -- Submitted, not yet accepted paper.
	\item \texttt{$\backslash$makepapertobesubmitted} -- Not yet submitted paper.
\end{itemize}

See code for this example document for examples on how to use.

\section{Literature Review... Cross-references}
Overview of major or important works for this topic
Cite
Find gap
Fill gap
%All labels throughout the thesis have to be unique. If the same
%equation or figure shows up twice e.g.\ a figure used
%both in the introduction and one of the included papers), it has
%to be given different labels. Referring to labels is done as
usual.

A simple trick to make sure this is the case and that will also help you keep track of all labels you used is to use the following naming convention:
\begin{itemize}
	\item \texttt{ch1:fig:labelname}, \texttt{ch1:tab:labelname}, \texttt{ch1:eq:labelname}, etc.\ all denote figures, tables and equations in Chapter 1.
	\item \texttt{paperA:fig:labelname}, \texttt{paperA:tab:labelname}, \texttt{paperA:eq:labelname}, etc.\ all denote figures, tables and equations in Paper A.
\end{itemize}

For existing text, e.g.\ papers, this is easily achieved by a simple search-and-replace operation on the string \texttt{$\backslash$label\{}. Any text editor will do that for you!

\section{Results Section.... Appendices\label{sec:app}}
What did I find?
What did I not find?
What I did find that was not expected?
%It is possible to have any number of appendices for each chapter.
%Simply type \texttt{$\backslash$appendix} before the first appendix,
%which is then a normal \texttt{section}, but numbered differently.

To add appendices to papers, use the
\texttt{$\backslash$paperappendix} command

\section{Dicussion Section... Including bibliography lists\label{sec:bib}}
Interpret results
Data supports goals?
Contribution or new?
Limitation (scope)
%In a thesis one might want several separate bibliographies. For
%example, one for the first part, and then separate bibliographies
%for each of the included papers.

This is solved using the \texttt{bibunits} package together with a
slight work-around in this template. For the first part of the
thesis, there is only one bibliography list, typeset like a chapter
(see this example document). In the papers, the bibliography lists
are typeset as un-numbered sections. See this file
\texttt{cseethesis\_example.tex} and \texttt{paper1.tex} for examples
how to place the bibliographies. 

Note that the command
\begin{itemize}
	\item \texttt{$\backslash$makebib} is used in Part I, to typeset the reference list in the thesis introduction.
	\item \texttt{$\backslash$putbib} is used in the papers in Part II.
\end{itemize}

\section{Conclusion Section... How to compile your project}
Finally, you probably like to know how to build your project to a
final PDF or PostScript file. Start by verifying that you can
compile this document. This is how it goes:
%
\begin{enumerate}
    \item Run \LaTeX\ (or pdfLaTex) once.
    \item Then run BibTeX on all the
    bu<i> files.
    \item Run \LaTeX\ twice more, to build the final DVI document
    (or pdfLaTeX if you want a PDF file).
\end{enumerate}

The above steps are easily collected in a script or batch file. See
the files \texttt{make.bat} and \texttt{compilebibunits.bat} for
examples.

\section{If I'm here is completed... Revision history}
By following :  \href{https://www.slideshare.net/NasreddineELGUEZAR/how-to-write-a-research-monograph-basics}{presentation} The template has evolved during several years and the exact
revisions are not clear to anyone. Starting from version 1.6,
however, the changes are more well-documented. This example
document will always support only the latest release of the
template. Below is a list of the revisions made to the document
class (and when applicable, the example document):%
\begin{itemize}
	 \item Version 3.1, September 1, 2010
		\begin{itemize}
			\item Fixed bug related to appendix numbering.
			\item Fixed page numbering of ``Part" pages.
			\item Added a comment at top of cseethesis\_example.tex, for improved compatibility with some editors.
		\end{itemize}
	 \item Version 3.0, June 7, 2009
		\begin{itemize}
			\item Fixed bug related to page headers in chapters containing no subsections.
		   \item Removed the EU class option and replaced with a logo argument to the preamble. See the code of this document for an example.
			\item The template is \textbf{no longer compatible with previous versions}.
         \item Removed the definition of boldface Greek letters from the document class, since this is not the proper place for that.
		\end{itemize}
  	 \item Version 2.5, March 5 2009
		\begin{itemize}
         \item Renamed the template \textit{cseethesis}. It is a continuation of the project initially called \textit{eisthesis}, but since its use has spread I decided to change the name.
			\item Various minor bug fixes.
			\item Update of example document, making examples of additions and revisions in recent versions of the template.
		\end{itemize}
	 \item Version 2.36: Added "Part" to the table of contents.
    \item Version 2.35:
        \begin{itemize}
				\item Fixed a bug regarding the page headers in the "appended papers part".
				\item Fixed a bug causing the section numbering to be wrong in a chapter succeeding a chapter containing appendices.
			   \item Added a chapter 3 in this example document, illustrating how to handle page headers for chapters without any sections. See the code at the top of \texttt{chapter3.tex}.
		  \end{itemize}
    \item Version 2.3:
        \begin{itemize}
            \item Added the commands \texttt{$\backslash$appendix}
            and \texttt{$\backslash$paperappendix}, see section
            \ref{sec:app}.
            \item The bibliography list is now typeset similar to a
            chapter in Part I, and as an un-numbered section in Part
            II. This required the use of separate commands for
            including the lists (see Sec.~\ref{sec:bib})
            \item Typesetting fixes for the table of contents page.
            As a consequence, the package \texttt{titletoc} is now
            required.
            \item Minor other code cleanup and bug fixes.
        \end{itemize}
    \item Version 2.25:
    \begin{itemize}
        \item Fixed a minor bug that used to generate a warning
        message regarding font shapes in the page headers.
    \end{itemize}
    \item Version 2.2:
    \begin{itemize}
        \item pdf and eps class options removed. The document class
        compiles with either pdfLaTeX or \LaTeX.
        \item The reference lists in papers are now typeset in the
        same way as in Part I of the thesis.
        \item Some minor adjustments of page header heights.
    \end{itemize}
    \item Version 2.1:
    \begin{itemize}
        \item Cross-references to chapters now work like they
        should. See the main document of this example.
        \item Major bug fixes to BibTeX reference lists, table of
        contents generation etc.
        \item The only change the user has to do is to use the new \texttt{$\backslash$makebib} instead of
        bibunits' \texttt{$\backslash$putbib}.
    \end{itemize}
    \item Version 2.0:
    \begin{itemize}
        \item Support for EU logotype on main page. The files
        \texttt{eu1\_f\_eng.pdf} and\\ \texttt{eu1\_f\_eng.eps} must be
        placed in the same directory as the document class.
        \item Support for pdf class options.
        \item Update of the \texttt{$\backslash$makechapter}
        command. It now requires three arguments. See main
        document for example.
        \item Support for BibTeX, using the \texttt{bibunits.sty}
        package.
    \end{itemize}
    \item Version 1.6: Various bug fixes to figure spacing etc.
\end{itemize}


\makechapter{Dicussion Chapter}{Table of contents}{Discussion of things\\ \label{ch5}}
\section{Discussion first section}
Interpret results
Data supports goals?
Contribution or new?
Limitation (scope)
%Current version of the \texttt{cseethesis} document class is: \classversion.\\
%Last modification: \classdate\\
%
%\noindent As of version 3.0, the template is no longer backwards compatible.
%
%\subsection{Sub1...About the document class}
%This document class was originally created in 2002 when I was working on my own PhD thesis. Since then, many people used it, found bugs (and occasionally even corrected them), and suggested improvements.
%
%The style is tailor-made for the typical types of theses that we write at the department, i.e.\ an introductory part followed by a collection of published or submitted research papers. 
%
%The template supports the use of both \LaTeX\ and pdfLaTeX. If you use the command \texttt{$\backslash$includegraphics} to import your figure and you supply the filename with its extension (e.g.\ .eps, .pdf), compilation should be possible with either one. For this to work, both .eps and .pdf versions of all figures must be available.
%
%The template is totally free to use, modify and distribute, as long as reference to the original author is kept and as long as all files remain in the package. Modified versions can only be distributed if it is clearly mentioned in the document class that modifications have been made and by whom.
%
%The whole package comes AS IS. I will correct bugs every now and then, but other than that, don't expect any support whatsoever.
%
%
%\subsection{Sub2...About this document}
%This document, as well as the actual \LaTeX\ code for it, makes up the documentation on how to use the document class.
%
%Only this chapter contains any readable information. Chapters 2 and 3 are only included as examples of some of the features of the template. The text is nonsense, but the corresponding \LaTeX\ code may be of some use. The same goes for the appended papers, which are only there as examples of a few options of the document class.
%
%Read this chapter carefully. If you have comments on what else should be in here in order to simplify the use of the document class, let me know.
%
%\section{Introduction...Chapters}
%\subsection{Defining chapters}
%Background info and context
%General focus
%Thesis statement (research Qs?)
%Attract attn.
%Why this topic
%%In order to add flexibility to the template, a new command, called \texttt{$\backslash$makechapter} is provided. This command takes three mandatory and one optional argument, as 
%\begin{center}
%	\texttt{$\backslash$makechapter[optional quote]\{page header\}\{toc entry\}\{Chapter title\}}
%\end{center}
%
%The reason this is solved like this is to allow for shorter page headers if the chapter name is very long. Also, if the actual chapter heading needs to be manually split in several lines (if the automatic splitting does not look so good), the table of contents (toc) entry might have to be defined differently. Note that normally, the last three arguments can be the same. 
%
%The use of an optional quote as an introduction to the chapter is demonstrated in this chapter. It can just as well be left out, which is demonstrated in this document (see the code).
%
%\subsection{Importing chapter contents}
%The sub-documents containing the chapters should start directly, i.e.\ they must not contain any $\backslash$begin\{document\} or $\backslash$end\{document\} tags.
%
%See this file, \textit{chapter1.tex} for details.
%
%\section{Methodology section...How to append papers}
%data collection analysis and interpretation how? Here!
%Description of the process
%
%%Papers are included using the \texttt{$\backslash$input} command, just as with chapters. You have to typeset paper title, authors, and abstract manually. See the example papers accompanying this document for an example.
%
%To make the separator sheet preceding each paper, use one of the following commands:
%\begin{itemize}
%	\item \texttt{$\backslash$makepaper} -- Published paper.
%	\item \texttt{$\backslash$makepaperaccepted} -- Accepted, not yet published paper.
%	\item \texttt{$\backslash$makepapersubmitted} -- Submitted, not yet accepted paper.
%	\item \texttt{$\backslash$makepapertobesubmitted} -- Not yet submitted paper.
%\end{itemize}
%
%See code for this example document for examples on how to use.
%
%\section{Literature Review... Cross-references}
%Overview of major or important works for this topic
%Cite
%Find gap
%Fill gap
%%All labels throughout the thesis have to be unique. If the same
%%equation or figure shows up twice e.g.\ a figure used
%%both in the introduction and one of the included papers), it has
%%to be given different labels. Referring to labels is done as
%usual.
%
%A simple trick to make sure this is the case and that will also help you keep track of all labels you used is to use the following naming convention:
%\begin{itemize}
%	\item \texttt{ch1:fig:labelname}, \texttt{ch1:tab:labelname}, \texttt{ch1:eq:labelname}, etc.\ all denote figures, tables and equations in Chapter 1.
%	\item \texttt{paperA:fig:labelname}, \texttt{paperA:tab:labelname}, \texttt{paperA:eq:labelname}, etc.\ all denote figures, tables and equations in Paper A.
%\end{itemize}
%
%For existing text, e.g.\ papers, this is easily achieved by a simple search-and-replace operation on the string \texttt{$\backslash$label\{}. Any text editor will do that for you!
%
%\section{Results Section.... Appendices\label{sec:app}}
%What did I find?
%What did I not find?
%What I did find that was not expected?
%%It is possible to have any number of appendices for each chapter.
%%Simply type \texttt{$\backslash$appendix} before the first appendix,
%%which is then a normal \texttt{section}, but numbered differently.
%
%To add appendices to papers, use the
%\texttt{$\backslash$paperappendix} command
%
%\section{Dicussion Section... Including bibliography lists\label{sec:bib}}
%Interpret results
%Data supports goals?
%Contribution or new?
%Limitation (scope)
%%In a thesis one might want several separate bibliographies. For
%%example, one for the first part, and then separate bibliographies
%%for each of the included papers.
%
%This is solved using the \texttt{bibunits} package together with a
%slight work-around in this template. For the first part of the
%thesis, there is only one bibliography list, typeset like a chapter
%(see this example document). In the papers, the bibliography lists
%are typeset as un-numbered sections. See this file
%\texttt{cseethesis\_example.tex} and \texttt{paper1.tex} for examples
%how to place the bibliographies. 
%
%Note that the command
%\begin{itemize}
%	\item \texttt{$\backslash$makebib} is used in Part I, to typeset the reference list in the thesis introduction.
%	\item \texttt{$\backslash$putbib} is used in the papers in Part II.
%\end{itemize}
%
%\section{Conclusion Section... How to compile your project}
%Finally, you probably like to know how to build your project to a
%final PDF or PostScript file. Start by verifying that you can
%compile this document. This is how it goes:
%%
%\begin{enumerate}
%    \item Run \LaTeX\ (or pdfLaTex) once.
%    \item Then run BibTeX on all the
%    bu<i> files.
%    \item Run \LaTeX\ twice more, to build the final DVI document
%    (or pdfLaTeX if you want a PDF file).
%\end{enumerate}
%
%The above steps are easily collected in a script or batch file. See
%the files \texttt{make.bat} and \texttt{compilebibunits.bat} for
%examples.
%
%\section{If I'm here is completed... Revision history}
%By following :  \href{https://www.slideshare.net/NasreddineELGUEZAR/how-to-write-a-research-monograph-basics}{presentation} The template has evolved during several years and the exact
%revisions are not clear to anyone. Starting from version 1.6,
%however, the changes are more well-documented. This example
%document will always support only the latest release of the
%template. Below is a list of the revisions made to the document
%class (and when applicable, the example document):%
%\begin{itemize}
%	 \item Version 3.1, September 1, 2010
%		\begin{itemize}
%			\item Fixed bug related to appendix numbering.
%			\item Fixed page numbering of ``Part" pages.
%			\item Added a comment at top of cseethesis\_example.tex, for improved compatibility with some editors.
%		\end{itemize}
%	 \item Version 3.0, June 7, 2009
%		\begin{itemize}
%			\item Fixed bug related to page headers in chapters containing no subsections.
%		   \item Removed the EU class option and replaced with a logo argument to the preamble. See the code of this document for an example.
%			\item The template is \textbf{no longer compatible with previous versions}.
%         \item Removed the definition of boldface Greek letters from the document class, since this is not the proper place for that.
%		\end{itemize}
%  	 \item Version 2.5, March 5 2009
%		\begin{itemize}
%         \item Renamed the template \textit{cseethesis}. It is a continuation of the project initially called \textit{eisthesis}, but since its use has spread I decided to change the name.
%			\item Various minor bug fixes.
%			\item Update of example document, making examples of additions and revisions in recent versions of the template.
%		\end{itemize}
%	 \item Version 2.36: Added "Part" to the table of contents.
%    \item Version 2.35:
%        \begin{itemize}
%				\item Fixed a bug regarding the page headers in the "appended papers part".
%				\item Fixed a bug causing the section numbering to be wrong in a chapter succeeding a chapter containing appendices.
%			   \item Added a chapter 3 in this example document, illustrating how to handle page headers for chapters without any sections. See the code at the top of \texttt{chapter3.tex}.
%		  \end{itemize}
%    \item Version 2.3:
%        \begin{itemize}
%            \item Added the commands \texttt{$\backslash$appendix}
%            and \texttt{$\backslash$paperappendix}, see section
%            \ref{sec:app}.
%            \item The bibliography list is now typeset similar to a
%            chapter in Part I, and as an un-numbered section in Part
%            II. This required the use of separate commands for
%            including the lists (see Sec.~\ref{sec:bib})
%            \item Typesetting fixes for the table of contents page.
%            As a consequence, the package \texttt{titletoc} is now
%            required.
%            \item Minor other code cleanup and bug fixes.
%        \end{itemize}
%    \item Version 2.25:
%    \begin{itemize}
%        \item Fixed a minor bug that used to generate a warning
%        message regarding font shapes in the page headers.
%    \end{itemize}
%    \item Version 2.2:
%    \begin{itemize}
%        \item pdf and eps class options removed. The document class
%        compiles with either pdfLaTeX or \LaTeX.
%        \item The reference lists in papers are now typeset in the
%        same way as in Part I of the thesis.
%        \item Some minor adjustments of page header heights.
%    \end{itemize}
%    \item Version 2.1:
%    \begin{itemize}
%        \item Cross-references to chapters now work like they
%        should. See the main document of this example.
%        \item Major bug fixes to BibTeX reference lists, table of
%        contents generation etc.
%        \item The only change the user has to do is to use the new \texttt{$\backslash$makebib} instead of
%        bibunits' \texttt{$\backslash$putbib}.
%    \end{itemize}
%    \item Version 2.0:
%    \begin{itemize}
%        \item Support for EU logotype on main page. The files
%        \texttt{eu1\_f\_eng.pdf} and\\ \texttt{eu1\_f\_eng.eps} must be
%        placed in the same directory as the document class.
%        \item Support for pdf class options.
%        \item Update of the \texttt{$\backslash$makechapter}
%        command. It now requires three arguments. See main
%        document for example.
%        \item Support for BibTeX, using the \texttt{bibunits.sty}
%        package.
%    \end{itemize}
%    \item Version 1.6: Various bug fixes to figure spacing etc.
%\end{itemize}


\makechapter{Conclusions  Chapter}{Table of contents}{Conclusions made, insights gathered\\ \label{ch6}}
\section{Conclusions first section}
Finally, you probably like to know how to build your project to a
final PDF or PostScript file. Start by verifying that you can
compile this document. This is how it goes:
Conclusions
\section{All in all where are we now}
After all this note what is being left to the reader and to my committee
%Current version of the \texttt{cseethesis} document class is: \classversion.\\
%Last modification: \classdate\\
%
%\noindent As of version 3.0, the template is no longer backwards compatible.
%
%\subsection{Sub1...About the document class}
%This document class was originally created in 2002 when I was working on my own PhD thesis. Since then, many people used it, found bugs (and occasionally even corrected them), and suggested improvements.
%
%The style is tailor-made for the typical types of theses that we write at the department, i.e.\ an introductory part followed by a collection of published or submitted research papers. 
%
%The template supports the use of both \LaTeX\ and pdfLaTeX. If you use the command \texttt{$\backslash$includegraphics} to import your figure and you supply the filename with its extension (e.g.\ .eps, .pdf), compilation should be possible with either one. For this to work, both .eps and .pdf versions of all figures must be available.
%
%The template is totally free to use, modify and distribute, as long as reference to the original author is kept and as long as all files remain in the package. Modified versions can only be distributed if it is clearly mentioned in the document class that modifications have been made and by whom.
%
%The whole package comes AS IS. I will correct bugs every now and then, but other than that, don't expect any support whatsoever.
%
%
%\subsection{Sub2...About this document}
%This document, as well as the actual \LaTeX\ code for it, makes up the documentation on how to use the document class.
%
%Only this chapter contains any readable information. Chapters 2 and 3 are only included as examples of some of the features of the template. The text is nonsense, but the corresponding \LaTeX\ code may be of some use. The same goes for the appended papers, which are only there as examples of a few options of the document class.
%
%Read this chapter carefully. If you have comments on what else should be in here in order to simplify the use of the document class, let me know.
%
%\section{Introduction...Chapters}
%\subsection{Defining chapters}
%Background info and context
%General focus
%Thesis statement (research Qs?)
%Attract attn.
%Why this topic
%%In order to add flexibility to the template, a new command, called \texttt{$\backslash$makechapter} is provided. This command takes three mandatory and one optional argument, as 
%\begin{center}
%	\texttt{$\backslash$makechapter[optional quote]\{page header\}\{toc entry\}\{Chapter title\}}
%\end{center}
%
%The reason this is solved like this is to allow for shorter page headers if the chapter name is very long. Also, if the actual chapter heading needs to be manually split in several lines (if the automatic splitting does not look so good), the table of contents (toc) entry might have to be defined differently. Note that normally, the last three arguments can be the same. 
%
%The use of an optional quote as an introduction to the chapter is demonstrated in this chapter. It can just as well be left out, which is demonstrated in this document (see the code).
%
%\subsection{Importing chapter contents}
%The sub-documents containing the chapters should start directly, i.e.\ they must not contain any $\backslash$begin\{document\} or $\backslash$end\{document\} tags.
%
%See this file, \textit{chapter1.tex} for details.
%
%\section{Methodology section...How to append papers}
%data collection analysis and interpretation how? Here!
%Description of the process
%
%%Papers are included using the \texttt{$\backslash$input} command, just as with chapters. You have to typeset paper title, authors, and abstract manually. See the example papers accompanying this document for an example.
%
%To make the separator sheet preceding each paper, use one of the following commands:
%\begin{itemize}
%	\item \texttt{$\backslash$makepaper} -- Published paper.
%	\item \texttt{$\backslash$makepaperaccepted} -- Accepted, not yet published paper.
%	\item \texttt{$\backslash$makepapersubmitted} -- Submitted, not yet accepted paper.
%	\item \texttt{$\backslash$makepapertobesubmitted} -- Not yet submitted paper.
%\end{itemize}
%
%See code for this example document for examples on how to use.
%
%\section{Literature Review... Cross-references}
%Overview of major or important works for this topic
%Cite
%Find gap
%Fill gap
%%All labels throughout the thesis have to be unique. If the same
%%equation or figure shows up twice e.g.\ a figure used
%%both in the introduction and one of the included papers), it has
%%to be given different labels. Referring to labels is done as
%usual.
%
%A simple trick to make sure this is the case and that will also help you keep track of all labels you used is to use the following naming convention:
%\begin{itemize}
%	\item \texttt{ch1:fig:labelname}, \texttt{ch1:tab:labelname}, \texttt{ch1:eq:labelname}, etc.\ all denote figures, tables and equations in Chapter 1.
%	\item \texttt{paperA:fig:labelname}, \texttt{paperA:tab:labelname}, \texttt{paperA:eq:labelname}, etc.\ all denote figures, tables and equations in Paper A.
%\end{itemize}
%
%For existing text, e.g.\ papers, this is easily achieved by a simple search-and-replace operation on the string \texttt{$\backslash$label\{}. Any text editor will do that for you!
%
%\section{Results Section.... Appendices\label{sec:app}}
%What did I find?
%What did I not find?
%What I did find that was not expected?
%%It is possible to have any number of appendices for each chapter.
%%Simply type \texttt{$\backslash$appendix} before the first appendix,
%%which is then a normal \texttt{section}, but numbered differently.
%
%To add appendices to papers, use the
%\texttt{$\backslash$paperappendix} command
%
%\section{Dicussion Section... Including bibliography lists\label{sec:bib}}
%Interpret results
%Data supports goals?
%Contribution or new?
%Limitation (scope)
%%In a thesis one might want several separate bibliographies. For
%%example, one for the first part, and then separate bibliographies
%%for each of the included papers.
%
%This is solved using the \texttt{bibunits} package together with a
%slight work-around in this template. For the first part of the
%thesis, there is only one bibliography list, typeset like a chapter
%(see this example document). In the papers, the bibliography lists
%are typeset as un-numbered sections. See this file
%\texttt{cseethesis\_example.tex} and \texttt{paper1.tex} for examples
%how to place the bibliographies. 
%
%Note that the command
%\begin{itemize}
%	\item \texttt{$\backslash$makebib} is used in Part I, to typeset the reference list in the thesis introduction.
%	\item \texttt{$\backslash$putbib} is used in the papers in Part II.
%\end{itemize}
%
%\section{Conclusion Section... How to compile your project}
%Finally, you probably like to know how to build your project to a
%final PDF or PostScript file. Start by verifying that you can
%compile this document. This is how it goes:
%%
%\begin{enumerate}
%    \item Run \LaTeX\ (or pdfLaTex) once.
%    \item Then run BibTeX on all the
%    bu<i> files.
%    \item Run \LaTeX\ twice more, to build the final DVI document
%    (or pdfLaTeX if you want a PDF file).
%\end{enumerate}
%
%The above steps are easily collected in a script or batch file. See
%the files \texttt{make.bat} and \texttt{compilebibunits.bat} for
%examples.
%
%\section{If I'm here is completed... Revision history}
%By following :  \href{https://www.slideshare.net/NasreddineELGUEZAR/how-to-write-a-research-monograph-basics}{presentation} The template has evolved during several years and the exact
%revisions are not clear to anyone. Starting from version 1.6,
%however, the changes are more well-documented. This example
%document will always support only the latest release of the
%template. Below is a list of the revisions made to the document
%class (and when applicable, the example document):%
%\begin{itemize}
%	 \item Version 3.1, September 1, 2010
%		\begin{itemize}
%			\item Fixed bug related to appendix numbering.
%			\item Fixed page numbering of ``Part" pages.
%			\item Added a comment at top of cseethesis\_example.tex, for improved compatibility with some editors.
%		\end{itemize}
%	 \item Version 3.0, June 7, 2009
%		\begin{itemize}
%			\item Fixed bug related to page headers in chapters containing no subsections.
%		   \item Removed the EU class option and replaced with a logo argument to the preamble. See the code of this document for an example.
%			\item The template is \textbf{no longer compatible with previous versions}.
%         \item Removed the definition of boldface Greek letters from the document class, since this is not the proper place for that.
%		\end{itemize}
%  	 \item Version 2.5, March 5 2009
%		\begin{itemize}
%         \item Renamed the template \textit{cseethesis}. It is a continuation of the project initially called \textit{eisthesis}, but since its use has spread I decided to change the name.
%			\item Various minor bug fixes.
%			\item Update of example document, making examples of additions and revisions in recent versions of the template.
%		\end{itemize}
%	 \item Version 2.36: Added "Part" to the table of contents.
%    \item Version 2.35:
%        \begin{itemize}
%				\item Fixed a bug regarding the page headers in the "appended papers part".
%				\item Fixed a bug causing the section numbering to be wrong in a chapter succeeding a chapter containing appendices.
%			   \item Added a chapter 3 in this example document, illustrating how to handle page headers for chapters without any sections. See the code at the top of \texttt{chapter3.tex}.
%		  \end{itemize}
%    \item Version 2.3:
%        \begin{itemize}
%            \item Added the commands \texttt{$\backslash$appendix}
%            and \texttt{$\backslash$paperappendix}, see section
%            \ref{sec:app}.
%            \item The bibliography list is now typeset similar to a
%            chapter in Part I, and as an un-numbered section in Part
%            II. This required the use of separate commands for
%            including the lists (see Sec.~\ref{sec:bib})
%            \item Typesetting fixes for the table of contents page.
%            As a consequence, the package \texttt{titletoc} is now
%            required.
%            \item Minor other code cleanup and bug fixes.
%        \end{itemize}
%    \item Version 2.25:
%    \begin{itemize}
%        \item Fixed a minor bug that used to generate a warning
%        message regarding font shapes in the page headers.
%    \end{itemize}
%    \item Version 2.2:
%    \begin{itemize}
%        \item pdf and eps class options removed. The document class
%        compiles with either pdfLaTeX or \LaTeX.
%        \item The reference lists in papers are now typeset in the
%        same way as in Part I of the thesis.
%        \item Some minor adjustments of page header heights.
%    \end{itemize}
%    \item Version 2.1:
%    \begin{itemize}
%        \item Cross-references to chapters now work like they
%        should. See the main document of this example.
%        \item Major bug fixes to BibTeX reference lists, table of
%        contents generation etc.
%        \item The only change the user has to do is to use the new \texttt{$\backslash$makebib} instead of
%        bibunits' \texttt{$\backslash$putbib}.
%    \end{itemize}
%    \item Version 2.0:
%    \begin{itemize}
%        \item Support for EU logotype on main page. The files
%        \texttt{eu1\_f\_eng.pdf} and\\ \texttt{eu1\_f\_eng.eps} must be
%        placed in the same directory as the document class.
%        \item Support for pdf class options.
%        \item Update of the \texttt{$\backslash$makechapter}
%        command. It now requires three arguments. See main
%        document for example.
%        \item Support for BibTeX, using the \texttt{bibunits.sty}
%        package.
%    \end{itemize}
%    \item Version 1.6: Various bug fixes to figure spacing etc.
%\end{itemize}


%\makechapter{Nonsense chapter}{Nonsense chapter}{Nonsense chapter, here only to verify that some issues in previous versions are really resolved!}
%This is a chapter with no sections, only here in order to test the document template. Please ignore the rest of this.
\vspace*{24pt}
dkhgfdsh fdsjhgiepy kdslkfds ewiuyfe fkjfdsohew dskjhfd fduew dsk di e sd lkjd dslkfd oiew sao cdkåwq acslkgå sdspo dsjdspe dkfo dkhgfdsh fdsjhgiepy kdslkfds ewiuyfe fkjfdsohew dskjhfd fduew dsk di e sd lkjd dslkfd oiew sao cdkåwq acslkgå sdspo dsjdspe dkfo dkhgfdsh fdsjhgiepy kdslkfds ewiuyfe fkjfdsohew dskjhfd fduew dsk di e sd lkjd dslkfd oiew sao cdkåwq acslkgå sdspo dsjdspe dkfo

dkhgfdsh fdsjhgiepy kdslkfds ewiuyfe fkjfdsohew dskjhfd fduew dsk di e sd lkjd dslkfd oiew sao cdkåwq acslkgå sdspo dsjdspe dkfodkhgfdsh fdsjhgiepy kdslkfds ewiuyfe fkjfdsohew dskjhfd fduew dsk di e sd lkjd dslkfd oiew sao cdkåwq acslkgå sdspo dsjdspe dkfo dkhgfdsh fdsjhgiepy kdslkfds ewiuyfe fkjfdsohew dskjhfd fduew dsk di e sd lkjd dslkfd oiew sao cdkåwq acslkgå sdspo dsjdspe dkfo dkhgfdsh fdsjhgiepy kdslkfds ewiuyfe fkjfdsohew dskjhfd fduew dsk di e sd lkjd dslkfd oiew sao cdkåwq acslkgå sdspo dsjdspe dkfo

dkhgfdsh fdsjhgiepy kdslkfds ewiuyfe fkjfdsohew dskjhfd fduew dsk di e sd lkjd dslkfd oiew sao cdkåwq acslkgå sdspo dsjdspe dkfodkhgfdsh fdsjhgiepy kdslkfds ewiuyfe fkjfdsohew dskjhfd fduew dsk di e sd lkjd dslkfd oiew sao cdkåwq acslkgå sdspo dsjdspe dkfo dkhgfdsh fdsjhgiepy kdslkfds ewiuyfe fkjfdsohew dskjhfd fduew dsk di e sd lkjd dslkfd oiew sao cdkåwq acslkgå sdspo dsjdspe dkfo dkhgfdsh fdsjhgiepy kdslkfds ewiuyfe fkjfdsohew dskjhfd fduew dsk di e sd lkjd dslkfd oiew sao cdkåwq acslkgå sdspo dsjdspe dkfo dkhgfdsh fdsjhgiepy kdslkfds ewiuyfe fkjfdsohew dskjhfd fduew dsk di e sd lkjd  dslkfd oiew sao cdkåwq acslkgå sdspo dsjdspe dkfodkhgfdsh fdsjhgiepy kdslkfds ewiuyfe fkjfdsohew dskjhfd fduew dsk di e sd lkjd dslkfd oiew sao cdkåwq acslkgå sdspo dsjdspe dkfo
dkhgfdsh fdsjhgiepy kdslkfds ewiuyfe fkjfdsohew dskjhfd fduew dsk di e sd lkjd dslkfd oiew sao cdkåwq acslkgå sdspo dsjdspe dkfo
dkhgfdsh fdsjhgiepy kdslkfds ewiuyfe fkjfdsohew dskjhfd fduew dsk di e sd lkjd dslkfd oiew sao cdkåwq acslkgå sdspo dsjdspe dkfo
dkhgfdsh fdsjhgiepy kdslkfds ewiuyfe fkjfdsohew dskjhfd fduew dsk di e sd lkjd dslkfd oiew sao cdkåwq acslkgå sdspo dsjdspe dkfo
dkhgfdsh fdsjhgiepy kdslkfds ewiuyfe fkjfdsohew dskjhfd fduew dsk di e sd lkjd dslkfd oiew sao cdkåwq acslkgå sdspo dsjdspe dkfo
dkhgfdsh fdsjhgiepy kdslkfds ewiuyfe fkjfdsohew dskjhfd fduew dsk di e sd lkjd dslkfd oiew sao cdkåwq acslkgå sdspo dsjdspe dkfo

dkhgfdsh fdsjhgiepy kdslkfds ewiuyfe fkjfdsohew dskjhfd fduew dsk di e sd lkjd dslkfd oiew sao cdkåwq acslkgå sdspo dsjdspe dkfo
dkhgfdsh fdsjhgiepy kdslkfds ewiuyfe fkjfdsohew dskjhfd fduew dsk di e sd lkjd dslkfd oiew sao cdkåwq acslkgå sdspo dsjdspe dkfo

dkhgfdsh fdsjhgiepy kdslkfds ewiuyfe fkjfdsohew dskjhfd fduew dsk di e sd lkjd dslkfd oiew sao cdkåwq acslkgå sdspo dsjdspe dkfo
 
dkhgfdsh fdsjhgiepy kdslkfds ewiuyfe fkjfdsohew dskjhfd fduew dsk di e sd lkjd dslkfd oiew sao cdkåwq acslkgå sdspo dsjdspe dkfodkhgfdsh fdsjhgiepy kdslkfds ewiuyfe fkjfdsohew dskjhfd fduew dsk di e sd lkjd dslkfd oiew sao cdkåwq acslkgå sdspo dsjdspe dkfodkhgfdsh fdsjhgiepy kdslkfds ewiuyfe fkjfdsohew dskjhfd fduew dsk di e sd lkjd dslkfd oiew sao cdkåwq acslkgå sdspo dsjdspe dkfodkhgfdsh fdsjhgiepy kdslkfds ewiuyfe fkjfdsohew dskjhfd fduew dsk di e sd lkjd dslkfd oiew sao cdkåwq acslkgå sdspo dsjdspe dkfodkhgfdsh fdsjhgiepy kdslkfds ewiuyfe fkjfdsohew dskjhfd fduew dsk di e sd lkjd dslkfd oiew sao cdkåwq acslkgå sdspo dsjdspe dkfodkhgfdsh fdsjhgiepy kdslkfds ewiuyfe fkjfdsohew dskjhfd fduew dsk di e sd lkjd dslkfd oiew sao cdkåwq acslkgå sdspo dsjdspe dkfodkhgfdsh fdsjhgiepy kdslkfds ewiuyfe fkjfdsohew dskjhfd fduew dsk di e sd lkjd dslkfd oiew sao cdkåwq acslkgå sdspo dsjdspe dkfodkhgfdsh fdsjhgiepy kdslkfds ewiuyfe fkjfdsohew dskjhfd fduew dsk di e sd lkjd dslkfd oiew sao cdkåwq acslkgå sdspo dsjdspe dkfodkhgfdsh fdsjhgiepy kdslkfds ewiuyfe fkjfdsohew dskjhfd fduew dsk di e sd lkjd dslkfd oiew sao cdkåwq acslkgå sdspo dsjdspe dkfodkhgfdsh fdsjhgiepy kdslkfds ewiuyfe fkjfdsohew dskjhfd fduew dsk di e sd lkjd dslkfd oiew sao cdkåwq acslkgå sdspo dsjdspe dkfodkhgfdsh fdsjhgiepy kdslkfds ewiuyfe fkjfdsohew dskjhfd fduew dsk di e sd lkjd dslkfd oiew sao cdkåwq acslkgå sdspo dsjdspe dkfo
dkhgfdsh fdsjhgiepy kdslkfds ewiuyfe fkjfdsohew dskjhfd fduew dsk di e sd lkjd dslkfd oiew sao cdkåwq acslkgå sdspo dsjdspe dkfodkhgfdsh fdsjhgiepy kdslkfds ewiuyfe fkjfdsohew dskjhfd fduew dsk di e sd lkjd dslkfd oiew sao cdkåwq acslkgå sdspo dsjdspe dkfo

dkhgfdsh fdsjhgiepy kdslkfds ewiuyfe fkjfdsohew dskjhfd fduew dsk di e sd lkjd dslkfd oiew sao cdkåwq acslkgå sdspo dsjdspe dkfodkhgfdsh fdsjhgiepy kdslkfds ewiuyfe fkjfdsohew dskjhfd fduew dsk di e sd lkjd dslkfd oiew sao cdkåwq acslkgå sdspo dsjdspe dkfodkhgfdsh fdsjhgiepy kdslkfds ewiuyfe fkjfdsohew dskjhfd fduew dsk di e sd lkjd dslkfd oiew sao cdkåwq acslkgå sdspo dsjdspe dkfodkhgfdsh fdsjhgiepy kdslkfds ewiuyfe fkjfdsohew dskjhfd fduew dsk di e sd lkjd dslkfd oiew sao cdkåwq acslkgå sdspo dsjdspe dkfo
dkhgfdsh fdsjhgiepy kdslkfds ewiuyfe fkjfdsohew dskjhfd fduew dsk di e sd lkjd dslkfd oiew sao cdkåwq acslkgå sdspo dsjdspe dkfo

dkhgfdsh fdsjhgiepy kdslkfds ewiuyfe fkjfdsohew dskjhfd fduew dsk di e sd lkjd dslkfd oiew sao cdkåwq acslkgå sdspo dsjdspe dkfodkhgfdsh fdsjhgiepy kdslkfds ewiuyfe fkjfdsohew dskjhfd fduew dsk di e sd lkjd dslkfd oiew sao cdkåwq acslkgå sdspo dsjdspe dkfodkhgfdsh fdsjhgiepy kdslkfds ewiuyfe fkjfdsohew dskjhfd fduew dsk di e sd lkjd dslkfd oiew sao cdkåwq acslkgå sdspo dsjdspe dkfodkhgfdsh fdsjhgiepy kdslkfds ewiuyfe fkjfdsohew dskjhfd fduew dsk di e sd lkjd dslkfd oiew sao cdkåwq acslkgå sdspo dsjdspe dkfodkhgfdsh fdsjhgiepy kdslkfds ewiuyfe fkjfdsohew dskjhfd fduew dsk di e sd lkjd dslkfd oiew sao cdkåwq acslkgå sdspo dsjdspe dkfo

dkhgfdsh fdsjhgiepy kdslkfds ewiuyfe fkjfdsohew dskjhfd fduew dsk di e sd lkjd dslkfd oiew sao cdkåwq acslkgå sdspo dsjdspe dkfodkhgfdsh fdsjhgiepy kdslkfds ewiuyfe fkjfdsohew dskjhfd fduew dsk di e sd lkjd dslkfd oiew sao cdkåwq acslkgå sdspo dsjdspe dkfodkhgfdsh fdsjhgiepy kdslkfds ewiuyfe fkjfdsohew dskjhfd fduew dsk di e sd lkjd dslkfd oiew sao cdkåwq acslkgå sdspo dsjdspe dkfodkhgfdsh fdsjhgiepy kdslkfds ewiuyfe fkjfdsohew dskjhfd fduew dsk di e sd lkjd dslkfd oiew sao cdkåwq acslkgå sdspo dsjdspe dkfodkhgfdsh fdsjhgiepy kdslkfds ewiuyfe fkjfdsohew dskjhfd fduew dsk di e sd lkjd dslkfd oiew sao cdkåwq acslkgå sdspo dsjdspe dkfodkhgfdsh fdsjhgiepy kdslkfds ewiuyfe fkjfdsohew dskjhfd fduew dsk di e sd lkjd dslkfd oiew sao cdkåwq acslkgå sdspo dsjdspe dkfodkhgfdsh fdsjhgiepy kdslkfds ewiuyfe fkjfdsohew dskjhfd fduew dsk di e sd lkjd dslkfd oiew sao cdkåwq acslkgå sdspo dsjdspe dkfo
dkhgfdsh fdsjhgiepy kdslkfds ewiuyfe fkjfdsohew dskjhfd fduew dsk di e sd lkjd dslkfd oiew sao cdkåwq acslkgå sdspo dsjdspe dkfodkhgfdsh fdsjhgiepy kdslkfds ewiuyfe fkjfdsohew dskjhfd fduew dsk di e sd lkjd dslkfd oiew sao cdkåwq acslkgå sdspo dsjdspe dkfo

dkhgfdsh fdsjhgiepy kdslkfds ewiuyfe fkjfdsohew dskjhfd fduew dsk di e sd lkjd dslkfd oiew sao cdkåwq acslkgå sdspo dsjdspe dkfodkhgfdsh fdsjhgiepy kdslkfds ewiuyfe fkjfdsohew dskjhfd fduew dsk di e sd lkjd dslkfd oiew sao cdkåwq acslkgå sdspo dsjdspe dkfodkhgfdsh fdsjhgiepy kdslkfds ewiuyfe fkjfdsohew dskjhfd fduew dsk di e sd lkjd dslkfd oiew sao cdkåwq acslkgå sdspo dsjdspe dkfodkhgfdsh fdsjhgiepy kdslkfds ewiuyfe fkjfdsohew dskjhfd fduew dsk di e sd lkjd dslkfd oiew sao cdkåwq acslkgå sdspo dsjdspe dkfodkhgfdsh fdsjhgiepy kdslkfds ewiuyfe fkjfdsohew dskjhfd fduew dsk di e sd lkjd dslkfd oiew sao cdkåwq acslkgå sdspo dsjdspe dkfo
dkhgfdsh fdsjhgiepy kdslkfds ewiuyfe fkjfdsohew dskjhfd fduew dsk di e sd lkjd dslkfd oiew sao cdkåwq acslkgå sdspo dsjdspe dkfodkhgfdsh fdsjhgiepy kdslkfds ewiuyfe fkjfdsohew dskjhfd fduew dsk di e sd lkjd dslkfd oiew sao cdkåwq acslkgå sdspo dsjdspe dkfo

dkhgfdsh fdsjhgiepy kdslkfds ewiuyfe fkjfdsohew dskjhfd fduew dsk di e sd lkjd dslkfd oiew sao cdkåwq acslkgå sdspo dsjdspe dkfodkhgfdsh fdsjhgiepy kdslkfds ewiuyfe fkjfdsohew dskjhfd fduew dsk di e sd lkjd dslkfd oiew sao cdkåwq acslkgå sdspo dsjdspe dkfodkhgfdsh fdsjhgiepy kdslkfds ewiuyfe fkjfdsohew dskjhfd fduew dsk di e sd lkjd dslkfd oiew sao cdkåwq acslkgå sdspo dsjdspe dkfo

dkhgfdsh fdsjhgiepy kdslkfds ewiuyfe fkjfdsohew dskjhfd fduew dsk di e sd lkjd dslkfd oiew sao cdkåwq acslkgå sdspo dsjdspe dkfodkhgfdsh fdsjhgiepy kdslkfds ewiuyfe fkjfdsohew dskjhfd fduew dsk di e sd lkjd dslkfd oiew sao cdkåwq acslkgå sdspo dsjdspe dkfodkhgfdsh fdsjhgiepy kdslkfds ewiuyfe fkjfdsohew dskjhfd fduew dsk di e sd lkjd dslkfd oiew sao cdkåwq acslkgå sdspo dsjdspe dkfodkhgfdsh fdsjhgiepy kdslkfds ewiuyfe fkjfdsohew dskjhfd fduew dsk di e sd lkjd dslkfd oiew sao cdkåwq acslkgå sdspo dsjdspe dkfodkhgfdsh fdsjhgiepy kdslkfds ewiuyfe fkjfdsohew dskjhfd fduew dsk di e sd lkjd dslkfd oiew sao cdkåwq acslkgå sdspo dsjdspe dkfodkhgfdsh fdsjhgiepy kdslkfds ewiuyfe fkjfdsohew dskjhfd fduew dsk di e sd lkjd dslkfd oiew sao cdkåwq acslkgå sdspo dsjdspe dkfodkhgfdsh fdsjhgiepy kdslkfds ewiuyfe fkjfdsohew dskjhfd fduew dsk di e sd lkjd dslkfd oiew sao cdkåwq acslkgå sdspo dsjdspe dkfo
dkhgfdsh fdsjhgiepy kdslkfds ewiuyfe fkjfdsohew dskjhfd fduew dsk di e sd lkjd dslkfd oiew sao cdkåwq acslkgå sdspo dsjdspe dkfodkhgfdsh fdsjhgiepy kdslkfds ewiuyfe fkjfdsohew dskjhfd fduew dsk di e sd lkjd dslkfd oiew sao cdkåwq acslkgå sdspo dsjdspe dkfo

dkhgfdsh fdsjhgiepy kdslkfds ewiuyfe fkjfdsohew dskjhfd fduew dsk di e sd lkjd dslkfd oiew sao cdkåwq acslkgå sdspo dsjdspe dkfodkhgfdsh fdsjhgiepy kdslkfds ewiuyfe fkjfdsohew dskjhfd fduew dsk di e sd lkjd dslkfd oiew sao cdkåwq acslkgå sdspo dsjdspe dkfodkhgfdsh fdsjhgiepy kdslkfds ewiuyfe fkjfdsohew dskjhfd fduew dsk di e sd lkjd dslkfd oiew sao cdkåwq acslkgå sdspo dsjdspe dkfodkhgfdsh fdsjhgiepy kdslkfds ewiuyfe fkjfdsohew dskjhfd fduew dsk di e sd lkjd dslkfd oiew sao cdkåwq acslkgå sdspo dsjdspe dkfo
dkhgfdsh fdsjhgiepy kdslkfds ewiuyfe fkjfdsohew dskjhfd fduew dsk di e sd lkjd dslkfd oiew sao cdkåwq acslkgå sdspo dsjdspe dkfodkhgfdsh fdsjhgiepy kdslkfds ewiuyfe fkjfdsohew dskjhfd fduew dsk di e sd lkjd dslkfd oiew sao cdkåwq acslkgå sdspo dsjdspe dkfodkhgfdsh fdsjhgiepy kdslkfds ewiuyfe fkjfdsohew dskjhfd fduew dsk di e sd lkjd dslkfd oiew sao cdkåwq acslkgå sdspo dsjdspe dkfodkhgfdsh fdsjhgiepy kdslkfds ewiuyfe fkjfdsohew dskjhfd fduew dsk di e sd lkjd dslkfd oiew sao cdkåwq acslkgå sdspo dsjdspe dkfo
dkhgfdsh fdsjhgiepy kdslkfds ewiuyfe fkjfdsohew dskjhfd fduew dsk di e sd lkjd dslkfd oiew sao cdkåwq acslkgå sdspo dsjdspe dkfodkhgfdsh fdsjhgiepy kdslkfds ewiuyfe fkjfdsohew dskjhfd fduew dsk di e sd lkjd dslkfd oiew sao cdkåwq acslkgå sdspo dsjdspe dkfodkhgfdsh fdsjhgiepy kdslkfds ewiuyfe fkjfdsohew dskjhfd fduew dsk di e sd lkjd dslkfd oiew sao cdkåwq acslkgå sdspo dsjdspe dkfodkhgfdsh fdsjhgiepy kdslkfds ewiuyfe fkjfdsohew dskjhfd fduew dsk di e sd lkjd dslkfd oiew sao cdkåwq acslkgå sdspo dsjdspe dkfo
dkhgfdsh fdsjhgiepy kdslkfds ewiuyfe fkjfdsohew dskjhfd fduew dsk di e sd lkjd dslkfd oiew sao cdkåwq acslkgå sdspo dsjdspe dkfodkhgfdsh fdsjhgiepy kdslkfds ewiuyfe fkjfdsohew dskjhfd fduew dsk di e sd lkjd dslkfd oiew sao cdkåwq acslkgå sdspo dsjdspe dkfodkhgfdsh fdsjhgiepy kdslkfds ewiuyfe fkjfdsohew dskjhfd fduew dsk di e sd lkjd dslkfd oiew sao cdkåwq acslkgå sdspo dsjdspe dkfo

dkhgfdsh fdsjhgiepy kdslkfds ewiuyfe fkjfdsohew dskjhfd fduew dsk di e sd lkjd dslkfd oiew sao cdkåwq acslkgå sdspo dsjdspe dkfodkhgfdsh fdsjhgiepy kdslkfds ewiuyfe fkjfdsohew dskjhfd fduew dsk di e sd lkjd dslkfd oiew sao cdkåwq acslkgå sdspo dsjdspe dkfodkhgfdsh fdsjhgiepy kdslkfds ewiuyfe fkjfdsohew dskjhfd fduew dsk di e sd lkjd dslkfd oiew sao cdkåwq acslkgå sdspo dsjdspe dkfodkhgfdsh fdsjhgiepy kdslkfds ewiuyfe fkjfdsohew dskjhfd fduew dsk di e sd lkjd dslkfd oiew sao cdkåwq acslkgå sdspo dsjdspe dkfodkhgfdsh fdsjhgiepy kdslkfds ewiuyfe fkjfdsohew dskjhfd fduew dsk di e sd lkjd dslkfd oiew sao cdkåwq acslkgå sdspo dsjdspe dkfodkhgfdsh fdsjhgiepy kdslkfds ewiuyfe fkjfdsohew dskjhfd fduew dsk di e sd lkjd dslkfd oiew sao cdkåwq acslkgå sdspo dsjdspe dkfodkhgfdsh fdsjhgiepy kdslkfds ewiuyfe fkjfdsohew dskjhfd fduew dsk di e sd lkjd dslkfd oiew sao cdkåwq acslkgå sdspo dsjdspe dkfodkhgfdsh fdsjhgiepy kdslkfds ewiuyfe fkjfdsohew dskjhfd fduew dsk di e sd lkjd dslkfd oiew sao cdkåwq acslkgå sdspo dsjdspe dkfodkhgfdsh fdsjhgiepy kdslkfds ewiuyfe fkjfdsohew dskjhfd fduew dsk di e sd lkjd dslkfd oiew sao cdkåwq acslkgå sdspo dsjdspe dkfodkhgfdsh fdsjhgiepy kdslkfds ewiuyfe fkjfdsohew dskjhfd fduew dsk di e sd lkjd dslkfd oiew sao cdkåwq acslkgå sdspo dsjdspe dkfodkhgfdsh fdsjhgiepy kdslkfds ewiuyfe fkjfdsohew dskjhfd fduew dsk di e sd lkjd dslkfd oiew sao cdkåwq acslkgå sdspo dsjdspe dkfo
dkhgfdsh fdsjhgiepy kdslkfds ewiuyfe fkjfdsohew dskjhfd fduew dsk di e sd lkjd dslkfd oiew sao cdkåwq acslkgå sdspo dsjdspe dkfodkhgfdsh fdsjhgiepy kdslkfds ewiuyfe fkjfdsohew dskjhfd fduew dsk di e sd lkjd dslkfd oiew sao cdkåwq acslkgå sdspo dsjdspe dkfo


%
%\makechapter{Nonsense chapter}{Nonsense chapter}{Nonsense chapter, here only to verify that some issues in previous versions are really resolved!}
%This is a chapter with no sections, only here in order to test the document template. Please ignore the rest of this.
\vspace*{24pt}
dkhgfdsh fdsjhgiepy kdslkfds ewiuyfe fkjfdsohew dskjhfd fduew dsk di e sd lkjd dslkfd oiew sao cdkåwq acslkgå sdspo dsjdspe dkfo dkhgfdsh fdsjhgiepy kdslkfds ewiuyfe fkjfdsohew dskjhfd fduew dsk di e sd lkjd dslkfd oiew sao cdkåwq acslkgå sdspo dsjdspe dkfo dkhgfdsh fdsjhgiepy kdslkfds ewiuyfe fkjfdsohew dskjhfd fduew dsk di e sd lkjd dslkfd oiew sao cdkåwq acslkgå sdspo dsjdspe dkfo

dkhgfdsh fdsjhgiepy kdslkfds ewiuyfe fkjfdsohew dskjhfd fduew dsk di e sd lkjd dslkfd oiew sao cdkåwq acslkgå sdspo dsjdspe dkfodkhgfdsh fdsjhgiepy kdslkfds ewiuyfe fkjfdsohew dskjhfd fduew dsk di e sd lkjd dslkfd oiew sao cdkåwq acslkgå sdspo dsjdspe dkfo dkhgfdsh fdsjhgiepy kdslkfds ewiuyfe fkjfdsohew dskjhfd fduew dsk di e sd lkjd dslkfd oiew sao cdkåwq acslkgå sdspo dsjdspe dkfo dkhgfdsh fdsjhgiepy kdslkfds ewiuyfe fkjfdsohew dskjhfd fduew dsk di e sd lkjd dslkfd oiew sao cdkåwq acslkgå sdspo dsjdspe dkfo

dkhgfdsh fdsjhgiepy kdslkfds ewiuyfe fkjfdsohew dskjhfd fduew dsk di e sd lkjd dslkfd oiew sao cdkåwq acslkgå sdspo dsjdspe dkfodkhgfdsh fdsjhgiepy kdslkfds ewiuyfe fkjfdsohew dskjhfd fduew dsk di e sd lkjd dslkfd oiew sao cdkåwq acslkgå sdspo dsjdspe dkfo dkhgfdsh fdsjhgiepy kdslkfds ewiuyfe fkjfdsohew dskjhfd fduew dsk di e sd lkjd dslkfd oiew sao cdkåwq acslkgå sdspo dsjdspe dkfo dkhgfdsh fdsjhgiepy kdslkfds ewiuyfe fkjfdsohew dskjhfd fduew dsk di e sd lkjd dslkfd oiew sao cdkåwq acslkgå sdspo dsjdspe dkfo dkhgfdsh fdsjhgiepy kdslkfds ewiuyfe fkjfdsohew dskjhfd fduew dsk di e sd lkjd  dslkfd oiew sao cdkåwq acslkgå sdspo dsjdspe dkfodkhgfdsh fdsjhgiepy kdslkfds ewiuyfe fkjfdsohew dskjhfd fduew dsk di e sd lkjd dslkfd oiew sao cdkåwq acslkgå sdspo dsjdspe dkfo
dkhgfdsh fdsjhgiepy kdslkfds ewiuyfe fkjfdsohew dskjhfd fduew dsk di e sd lkjd dslkfd oiew sao cdkåwq acslkgå sdspo dsjdspe dkfo
dkhgfdsh fdsjhgiepy kdslkfds ewiuyfe fkjfdsohew dskjhfd fduew dsk di e sd lkjd dslkfd oiew sao cdkåwq acslkgå sdspo dsjdspe dkfo
dkhgfdsh fdsjhgiepy kdslkfds ewiuyfe fkjfdsohew dskjhfd fduew dsk di e sd lkjd dslkfd oiew sao cdkåwq acslkgå sdspo dsjdspe dkfo
dkhgfdsh fdsjhgiepy kdslkfds ewiuyfe fkjfdsohew dskjhfd fduew dsk di e sd lkjd dslkfd oiew sao cdkåwq acslkgå sdspo dsjdspe dkfo
dkhgfdsh fdsjhgiepy kdslkfds ewiuyfe fkjfdsohew dskjhfd fduew dsk di e sd lkjd dslkfd oiew sao cdkåwq acslkgå sdspo dsjdspe dkfo

dkhgfdsh fdsjhgiepy kdslkfds ewiuyfe fkjfdsohew dskjhfd fduew dsk di e sd lkjd dslkfd oiew sao cdkåwq acslkgå sdspo dsjdspe dkfo
dkhgfdsh fdsjhgiepy kdslkfds ewiuyfe fkjfdsohew dskjhfd fduew dsk di e sd lkjd dslkfd oiew sao cdkåwq acslkgå sdspo dsjdspe dkfo

dkhgfdsh fdsjhgiepy kdslkfds ewiuyfe fkjfdsohew dskjhfd fduew dsk di e sd lkjd dslkfd oiew sao cdkåwq acslkgå sdspo dsjdspe dkfo
 
dkhgfdsh fdsjhgiepy kdslkfds ewiuyfe fkjfdsohew dskjhfd fduew dsk di e sd lkjd dslkfd oiew sao cdkåwq acslkgå sdspo dsjdspe dkfodkhgfdsh fdsjhgiepy kdslkfds ewiuyfe fkjfdsohew dskjhfd fduew dsk di e sd lkjd dslkfd oiew sao cdkåwq acslkgå sdspo dsjdspe dkfodkhgfdsh fdsjhgiepy kdslkfds ewiuyfe fkjfdsohew dskjhfd fduew dsk di e sd lkjd dslkfd oiew sao cdkåwq acslkgå sdspo dsjdspe dkfodkhgfdsh fdsjhgiepy kdslkfds ewiuyfe fkjfdsohew dskjhfd fduew dsk di e sd lkjd dslkfd oiew sao cdkåwq acslkgå sdspo dsjdspe dkfodkhgfdsh fdsjhgiepy kdslkfds ewiuyfe fkjfdsohew dskjhfd fduew dsk di e sd lkjd dslkfd oiew sao cdkåwq acslkgå sdspo dsjdspe dkfodkhgfdsh fdsjhgiepy kdslkfds ewiuyfe fkjfdsohew dskjhfd fduew dsk di e sd lkjd dslkfd oiew sao cdkåwq acslkgå sdspo dsjdspe dkfodkhgfdsh fdsjhgiepy kdslkfds ewiuyfe fkjfdsohew dskjhfd fduew dsk di e sd lkjd dslkfd oiew sao cdkåwq acslkgå sdspo dsjdspe dkfodkhgfdsh fdsjhgiepy kdslkfds ewiuyfe fkjfdsohew dskjhfd fduew dsk di e sd lkjd dslkfd oiew sao cdkåwq acslkgå sdspo dsjdspe dkfodkhgfdsh fdsjhgiepy kdslkfds ewiuyfe fkjfdsohew dskjhfd fduew dsk di e sd lkjd dslkfd oiew sao cdkåwq acslkgå sdspo dsjdspe dkfodkhgfdsh fdsjhgiepy kdslkfds ewiuyfe fkjfdsohew dskjhfd fduew dsk di e sd lkjd dslkfd oiew sao cdkåwq acslkgå sdspo dsjdspe dkfodkhgfdsh fdsjhgiepy kdslkfds ewiuyfe fkjfdsohew dskjhfd fduew dsk di e sd lkjd dslkfd oiew sao cdkåwq acslkgå sdspo dsjdspe dkfo
dkhgfdsh fdsjhgiepy kdslkfds ewiuyfe fkjfdsohew dskjhfd fduew dsk di e sd lkjd dslkfd oiew sao cdkåwq acslkgå sdspo dsjdspe dkfodkhgfdsh fdsjhgiepy kdslkfds ewiuyfe fkjfdsohew dskjhfd fduew dsk di e sd lkjd dslkfd oiew sao cdkåwq acslkgå sdspo dsjdspe dkfo

dkhgfdsh fdsjhgiepy kdslkfds ewiuyfe fkjfdsohew dskjhfd fduew dsk di e sd lkjd dslkfd oiew sao cdkåwq acslkgå sdspo dsjdspe dkfodkhgfdsh fdsjhgiepy kdslkfds ewiuyfe fkjfdsohew dskjhfd fduew dsk di e sd lkjd dslkfd oiew sao cdkåwq acslkgå sdspo dsjdspe dkfodkhgfdsh fdsjhgiepy kdslkfds ewiuyfe fkjfdsohew dskjhfd fduew dsk di e sd lkjd dslkfd oiew sao cdkåwq acslkgå sdspo dsjdspe dkfodkhgfdsh fdsjhgiepy kdslkfds ewiuyfe fkjfdsohew dskjhfd fduew dsk di e sd lkjd dslkfd oiew sao cdkåwq acslkgå sdspo dsjdspe dkfo
dkhgfdsh fdsjhgiepy kdslkfds ewiuyfe fkjfdsohew dskjhfd fduew dsk di e sd lkjd dslkfd oiew sao cdkåwq acslkgå sdspo dsjdspe dkfo

dkhgfdsh fdsjhgiepy kdslkfds ewiuyfe fkjfdsohew dskjhfd fduew dsk di e sd lkjd dslkfd oiew sao cdkåwq acslkgå sdspo dsjdspe dkfodkhgfdsh fdsjhgiepy kdslkfds ewiuyfe fkjfdsohew dskjhfd fduew dsk di e sd lkjd dslkfd oiew sao cdkåwq acslkgå sdspo dsjdspe dkfodkhgfdsh fdsjhgiepy kdslkfds ewiuyfe fkjfdsohew dskjhfd fduew dsk di e sd lkjd dslkfd oiew sao cdkåwq acslkgå sdspo dsjdspe dkfodkhgfdsh fdsjhgiepy kdslkfds ewiuyfe fkjfdsohew dskjhfd fduew dsk di e sd lkjd dslkfd oiew sao cdkåwq acslkgå sdspo dsjdspe dkfodkhgfdsh fdsjhgiepy kdslkfds ewiuyfe fkjfdsohew dskjhfd fduew dsk di e sd lkjd dslkfd oiew sao cdkåwq acslkgå sdspo dsjdspe dkfo

dkhgfdsh fdsjhgiepy kdslkfds ewiuyfe fkjfdsohew dskjhfd fduew dsk di e sd lkjd dslkfd oiew sao cdkåwq acslkgå sdspo dsjdspe dkfodkhgfdsh fdsjhgiepy kdslkfds ewiuyfe fkjfdsohew dskjhfd fduew dsk di e sd lkjd dslkfd oiew sao cdkåwq acslkgå sdspo dsjdspe dkfodkhgfdsh fdsjhgiepy kdslkfds ewiuyfe fkjfdsohew dskjhfd fduew dsk di e sd lkjd dslkfd oiew sao cdkåwq acslkgå sdspo dsjdspe dkfodkhgfdsh fdsjhgiepy kdslkfds ewiuyfe fkjfdsohew dskjhfd fduew dsk di e sd lkjd dslkfd oiew sao cdkåwq acslkgå sdspo dsjdspe dkfodkhgfdsh fdsjhgiepy kdslkfds ewiuyfe fkjfdsohew dskjhfd fduew dsk di e sd lkjd dslkfd oiew sao cdkåwq acslkgå sdspo dsjdspe dkfodkhgfdsh fdsjhgiepy kdslkfds ewiuyfe fkjfdsohew dskjhfd fduew dsk di e sd lkjd dslkfd oiew sao cdkåwq acslkgå sdspo dsjdspe dkfodkhgfdsh fdsjhgiepy kdslkfds ewiuyfe fkjfdsohew dskjhfd fduew dsk di e sd lkjd dslkfd oiew sao cdkåwq acslkgå sdspo dsjdspe dkfo
dkhgfdsh fdsjhgiepy kdslkfds ewiuyfe fkjfdsohew dskjhfd fduew dsk di e sd lkjd dslkfd oiew sao cdkåwq acslkgå sdspo dsjdspe dkfodkhgfdsh fdsjhgiepy kdslkfds ewiuyfe fkjfdsohew dskjhfd fduew dsk di e sd lkjd dslkfd oiew sao cdkåwq acslkgå sdspo dsjdspe dkfo

dkhgfdsh fdsjhgiepy kdslkfds ewiuyfe fkjfdsohew dskjhfd fduew dsk di e sd lkjd dslkfd oiew sao cdkåwq acslkgå sdspo dsjdspe dkfodkhgfdsh fdsjhgiepy kdslkfds ewiuyfe fkjfdsohew dskjhfd fduew dsk di e sd lkjd dslkfd oiew sao cdkåwq acslkgå sdspo dsjdspe dkfodkhgfdsh fdsjhgiepy kdslkfds ewiuyfe fkjfdsohew dskjhfd fduew dsk di e sd lkjd dslkfd oiew sao cdkåwq acslkgå sdspo dsjdspe dkfodkhgfdsh fdsjhgiepy kdslkfds ewiuyfe fkjfdsohew dskjhfd fduew dsk di e sd lkjd dslkfd oiew sao cdkåwq acslkgå sdspo dsjdspe dkfodkhgfdsh fdsjhgiepy kdslkfds ewiuyfe fkjfdsohew dskjhfd fduew dsk di e sd lkjd dslkfd oiew sao cdkåwq acslkgå sdspo dsjdspe dkfo
dkhgfdsh fdsjhgiepy kdslkfds ewiuyfe fkjfdsohew dskjhfd fduew dsk di e sd lkjd dslkfd oiew sao cdkåwq acslkgå sdspo dsjdspe dkfodkhgfdsh fdsjhgiepy kdslkfds ewiuyfe fkjfdsohew dskjhfd fduew dsk di e sd lkjd dslkfd oiew sao cdkåwq acslkgå sdspo dsjdspe dkfo

dkhgfdsh fdsjhgiepy kdslkfds ewiuyfe fkjfdsohew dskjhfd fduew dsk di e sd lkjd dslkfd oiew sao cdkåwq acslkgå sdspo dsjdspe dkfodkhgfdsh fdsjhgiepy kdslkfds ewiuyfe fkjfdsohew dskjhfd fduew dsk di e sd lkjd dslkfd oiew sao cdkåwq acslkgå sdspo dsjdspe dkfodkhgfdsh fdsjhgiepy kdslkfds ewiuyfe fkjfdsohew dskjhfd fduew dsk di e sd lkjd dslkfd oiew sao cdkåwq acslkgå sdspo dsjdspe dkfo

dkhgfdsh fdsjhgiepy kdslkfds ewiuyfe fkjfdsohew dskjhfd fduew dsk di e sd lkjd dslkfd oiew sao cdkåwq acslkgå sdspo dsjdspe dkfodkhgfdsh fdsjhgiepy kdslkfds ewiuyfe fkjfdsohew dskjhfd fduew dsk di e sd lkjd dslkfd oiew sao cdkåwq acslkgå sdspo dsjdspe dkfodkhgfdsh fdsjhgiepy kdslkfds ewiuyfe fkjfdsohew dskjhfd fduew dsk di e sd lkjd dslkfd oiew sao cdkåwq acslkgå sdspo dsjdspe dkfodkhgfdsh fdsjhgiepy kdslkfds ewiuyfe fkjfdsohew dskjhfd fduew dsk di e sd lkjd dslkfd oiew sao cdkåwq acslkgå sdspo dsjdspe dkfodkhgfdsh fdsjhgiepy kdslkfds ewiuyfe fkjfdsohew dskjhfd fduew dsk di e sd lkjd dslkfd oiew sao cdkåwq acslkgå sdspo dsjdspe dkfodkhgfdsh fdsjhgiepy kdslkfds ewiuyfe fkjfdsohew dskjhfd fduew dsk di e sd lkjd dslkfd oiew sao cdkåwq acslkgå sdspo dsjdspe dkfodkhgfdsh fdsjhgiepy kdslkfds ewiuyfe fkjfdsohew dskjhfd fduew dsk di e sd lkjd dslkfd oiew sao cdkåwq acslkgå sdspo dsjdspe dkfo
dkhgfdsh fdsjhgiepy kdslkfds ewiuyfe fkjfdsohew dskjhfd fduew dsk di e sd lkjd dslkfd oiew sao cdkåwq acslkgå sdspo dsjdspe dkfodkhgfdsh fdsjhgiepy kdslkfds ewiuyfe fkjfdsohew dskjhfd fduew dsk di e sd lkjd dslkfd oiew sao cdkåwq acslkgå sdspo dsjdspe dkfo

dkhgfdsh fdsjhgiepy kdslkfds ewiuyfe fkjfdsohew dskjhfd fduew dsk di e sd lkjd dslkfd oiew sao cdkåwq acslkgå sdspo dsjdspe dkfodkhgfdsh fdsjhgiepy kdslkfds ewiuyfe fkjfdsohew dskjhfd fduew dsk di e sd lkjd dslkfd oiew sao cdkåwq acslkgå sdspo dsjdspe dkfodkhgfdsh fdsjhgiepy kdslkfds ewiuyfe fkjfdsohew dskjhfd fduew dsk di e sd lkjd dslkfd oiew sao cdkåwq acslkgå sdspo dsjdspe dkfodkhgfdsh fdsjhgiepy kdslkfds ewiuyfe fkjfdsohew dskjhfd fduew dsk di e sd lkjd dslkfd oiew sao cdkåwq acslkgå sdspo dsjdspe dkfo
dkhgfdsh fdsjhgiepy kdslkfds ewiuyfe fkjfdsohew dskjhfd fduew dsk di e sd lkjd dslkfd oiew sao cdkåwq acslkgå sdspo dsjdspe dkfodkhgfdsh fdsjhgiepy kdslkfds ewiuyfe fkjfdsohew dskjhfd fduew dsk di e sd lkjd dslkfd oiew sao cdkåwq acslkgå sdspo dsjdspe dkfodkhgfdsh fdsjhgiepy kdslkfds ewiuyfe fkjfdsohew dskjhfd fduew dsk di e sd lkjd dslkfd oiew sao cdkåwq acslkgå sdspo dsjdspe dkfodkhgfdsh fdsjhgiepy kdslkfds ewiuyfe fkjfdsohew dskjhfd fduew dsk di e sd lkjd dslkfd oiew sao cdkåwq acslkgå sdspo dsjdspe dkfo
dkhgfdsh fdsjhgiepy kdslkfds ewiuyfe fkjfdsohew dskjhfd fduew dsk di e sd lkjd dslkfd oiew sao cdkåwq acslkgå sdspo dsjdspe dkfodkhgfdsh fdsjhgiepy kdslkfds ewiuyfe fkjfdsohew dskjhfd fduew dsk di e sd lkjd dslkfd oiew sao cdkåwq acslkgå sdspo dsjdspe dkfodkhgfdsh fdsjhgiepy kdslkfds ewiuyfe fkjfdsohew dskjhfd fduew dsk di e sd lkjd dslkfd oiew sao cdkåwq acslkgå sdspo dsjdspe dkfodkhgfdsh fdsjhgiepy kdslkfds ewiuyfe fkjfdsohew dskjhfd fduew dsk di e sd lkjd dslkfd oiew sao cdkåwq acslkgå sdspo dsjdspe dkfo
dkhgfdsh fdsjhgiepy kdslkfds ewiuyfe fkjfdsohew dskjhfd fduew dsk di e sd lkjd dslkfd oiew sao cdkåwq acslkgå sdspo dsjdspe dkfodkhgfdsh fdsjhgiepy kdslkfds ewiuyfe fkjfdsohew dskjhfd fduew dsk di e sd lkjd dslkfd oiew sao cdkåwq acslkgå sdspo dsjdspe dkfodkhgfdsh fdsjhgiepy kdslkfds ewiuyfe fkjfdsohew dskjhfd fduew dsk di e sd lkjd dslkfd oiew sao cdkåwq acslkgå sdspo dsjdspe dkfo

dkhgfdsh fdsjhgiepy kdslkfds ewiuyfe fkjfdsohew dskjhfd fduew dsk di e sd lkjd dslkfd oiew sao cdkåwq acslkgå sdspo dsjdspe dkfodkhgfdsh fdsjhgiepy kdslkfds ewiuyfe fkjfdsohew dskjhfd fduew dsk di e sd lkjd dslkfd oiew sao cdkåwq acslkgå sdspo dsjdspe dkfodkhgfdsh fdsjhgiepy kdslkfds ewiuyfe fkjfdsohew dskjhfd fduew dsk di e sd lkjd dslkfd oiew sao cdkåwq acslkgå sdspo dsjdspe dkfodkhgfdsh fdsjhgiepy kdslkfds ewiuyfe fkjfdsohew dskjhfd fduew dsk di e sd lkjd dslkfd oiew sao cdkåwq acslkgå sdspo dsjdspe dkfodkhgfdsh fdsjhgiepy kdslkfds ewiuyfe fkjfdsohew dskjhfd fduew dsk di e sd lkjd dslkfd oiew sao cdkåwq acslkgå sdspo dsjdspe dkfodkhgfdsh fdsjhgiepy kdslkfds ewiuyfe fkjfdsohew dskjhfd fduew dsk di e sd lkjd dslkfd oiew sao cdkåwq acslkgå sdspo dsjdspe dkfodkhgfdsh fdsjhgiepy kdslkfds ewiuyfe fkjfdsohew dskjhfd fduew dsk di e sd lkjd dslkfd oiew sao cdkåwq acslkgå sdspo dsjdspe dkfodkhgfdsh fdsjhgiepy kdslkfds ewiuyfe fkjfdsohew dskjhfd fduew dsk di e sd lkjd dslkfd oiew sao cdkåwq acslkgå sdspo dsjdspe dkfodkhgfdsh fdsjhgiepy kdslkfds ewiuyfe fkjfdsohew dskjhfd fduew dsk di e sd lkjd dslkfd oiew sao cdkåwq acslkgå sdspo dsjdspe dkfodkhgfdsh fdsjhgiepy kdslkfds ewiuyfe fkjfdsohew dskjhfd fduew dsk di e sd lkjd dslkfd oiew sao cdkåwq acslkgå sdspo dsjdspe dkfodkhgfdsh fdsjhgiepy kdslkfds ewiuyfe fkjfdsohew dskjhfd fduew dsk di e sd lkjd dslkfd oiew sao cdkåwq acslkgå sdspo dsjdspe dkfo
dkhgfdsh fdsjhgiepy kdslkfds ewiuyfe fkjfdsohew dskjhfd fduew dsk di e sd lkjd dslkfd oiew sao cdkåwq acslkgå sdspo dsjdspe dkfodkhgfdsh fdsjhgiepy kdslkfds ewiuyfe fkjfdsohew dskjhfd fduew dsk di e sd lkjd dslkfd oiew sao cdkåwq acslkgå sdspo dsjdspe dkfo



\makebib
\end{bibunit}

%%%%%%%%%%%%%%%%%%%%%%%%%%%%%%%%%%%%%%%%%%%%%%%%%%%%%%%%%%%%%%%%%%%%
%% Begin Part II - Collection of papers
%%%%%%%%%%%%%%%%%%%%%%%%%%%%%%%%%%%%%%%%%%%%%%%%%%%%%%%%%%%%%%%%%%%%

%\makepartpage{Part II}%
%\startpapers
%
%%-------------------------------------------------------------------
%\def\paperheader{Paper A}
%\def\papertitle{The Title of the Papers in the Thesis are Automatically Split In Several Lines if Necessary}
%\def\paperauthorstring{John Doe and Jane Doe}
%\def\referencestring{Example Thesis, Internal Report, Lule{\aa} University of Technology, 2002.}
%\def\copyrightstring{2002, The Publisher, Reprinted with permission.}
%
%% The definitions above could just as well be put directly into the function
%% call below, but were explicitly defined to more clearly illustrate the
%% use of the function \makepaper.
%
%\makepaper
%  {\paperheader}
%  {\papertitle}
%  {\paperauthorstring}
%  {\referencestring}
%  {\copyrightstring}
%
%% The actual contents is imported by un-commenting the \input line below.
%% Make sure the file exist.
%\begin{bibunit}
\thispagestyle{plain}
\begin{center}
\Large\textbf{The Tile of the Paper}\\
\vspace{10mm}
\normalsize John Doe and Jane Doe\\
\vspace{15mm}
\textbf{Abstract}\\
\end{center}
Abstract text of the paper...

\section{Introduction}
The text of this article is imported from the file
\textit{paper1.tex}. The title and abstract part above are typeset
manually (see file for code template).
\cite{Goossens}

\acrfull{nlp} is the full version of this short version \acrshort{nlp} and the long version \acrlong{nlp}. 

Be sure NOT to have any $\backslash$begin\{document\} or
$\backslash$end\{document\} tags in the imported files.

Some references here too, just to show the use of bibunits
\nocite{*}.

\paperappendix
\section{First appendix of paper A}
Some appendix text.
\subsection{A subsection of the appendix}
%
\begin{equation}
    X(\omega) = \int_{-\infty}^{\infty} x(t)e^{-j\omega t}\,dt.
\end{equation}
%
\subsection{Another subsection of the appendix}
\subsubsection{Test subsubsection}

\section{Another appendix}
Some text in the second appendix

%%% Put references here
\putbib
\end{bibunit}

%
%
%%-------------------------------------------------------------------
%\def\paperheader{Paper B}
%\def\papertitle{The Theory of Research}
%\def\paperauthorstring{John Doe and Jane Doe}
%\def\referencestring{Example Thesis, Internal Report, Lule{\aa} University of Technology, 2009.}
%\def\copyrightstring{2009, The Publisher, Reprinted with permission.}
%
%% The definitions above could just as well be put directly into the function
%% call below, but were explicitly defined to more clearly illustrate the
%% use of the function \makepaper.
%
%\makepaperaccepted
%  {\paperheader}
%  {\papertitle}
%  {\paperauthorstring}
%  {\referencestring}
%  {\copyrightstring}
%
%% The actual contents is imported by un-commenting the \input line below.
%% Make sure the file exist.
%\begin{bibunit}
\thispagestyle{plain}
\begin{center}
\Large\textbf{\papertitle}\\
\vspace{10mm}
\normalsize John Doe and Jane Doe\\
\vspace{15mm}
\textbf{Abstract}\\
\end{center}
Abstract text of the paper...

\section{Introduction}
The text of this article is imported from the file
\textit{paper2.tex}. The title and abstract part above are typeset
manually (see file for code template).

Be sure NOT to have any $\backslash$begin\{document\} or
$\backslash$end\{document\} tags in the imported files.

Some references here too, just to show the use of bibunits
\nocite{*}.

\paperappendix
\section{First appendix of paper B}
Some appendix text.
\subsection{A subsection of the appendix}
%
\begin{equation}
    X(\omega) = \int_{-\infty}^{\infty} x(t)e^{-j\omega t}\,dt.
\end{equation}
%
\subsection{Another subsection of the appendix}

\section{Another appendix}
Some text in the second appendix

%%% Put references here
\putbib
\end{bibunit}

%
%%-------------------------------------------------------------------
%\def\paperheader{Paper C}
%\def\papertitle{Yet Another Sub-Optimal Estimator of Sinusoids in Noise}
%\def\paperauthorstring{Dr.\ C}
%\def\referencestring{Example Thesis, Internal Report, Lule{\aa} University of Technology, 2009.}
%\def\copyrightstring{2009, The Publisher, Reprinted with permission.}
%
%% The definitions above could just as well be put directly into the function
%% call below, but were explicitly defined to more clearly illustrate the
%% use of the function \makepaper.
%
%\makepapersubmitted
%  {\paperheader}
%  {\papertitle}
%  {\paperauthorstring}
%  {\referencestring}
%  {\copyrightstring}
%
%% The actual contents is imported by un-commenting the \input line below.
%% Make sure the file exist.
%\begin{bibunit}
\thispagestyle{plain}
\begin{center}
\Large\textbf{\papertitle}\\
\vspace{10mm}
\normalsize Dr.~C\\
\vspace{15mm}
\textbf{Abstract}\\
\end{center}
Abstract text of the paper...

\section{Introduction}
The text of this article is imported from the file
\textit{paper3.tex}. The title and abstract part above are typeset
manually (see file for code template).

Be sure NOT to have any $\backslash$begin\{document\} or
$\backslash$end\{document\} tags in the imported files.

Some references here too, just to show the use of bibunits
\nocite{*}.


%%% Put references here
\putbib
\end{bibunit}

%
%%-------------------------------------------------------------------
%\def\paperheader{Paper D}
%\def\papertitle{An example of a yet-to-be-submitted paper}
%\def\paperauthorstring{Dr.\ C}
%
%% The definitions above could just as well be put directly into the function
%% call below, but were explicitly defined to more clearly illustrate the
%% use of the function \makepaper.
%
%\makepapertobesubmitted
%  {\paperheader}
%  {\papertitle}
%  {\paperauthorstring}
%
%% The actual contents is imported by un-commenting the \input line below.
%% Make sure the file exist.
%\begin{bibunit}
\thispagestyle{plain}
\begin{center}
\Large\textbf{\papertitle}\\
\vspace{10mm}
\normalsize Dr.~C\\
\vspace{15mm}
\textbf{Abstract}\\
\end{center}
Abstract text of the paper...

\section{Introduction}
The text of this article is imported from the file
\textit{paper3.tex}. The title and abstract part above are typeset
manually (see file for code template).

Be sure NOT to have any $\backslash$begin\{document\} or
$\backslash$end\{document\} tags in the imported files.

Some references here too, just to show the use of bibunits
\nocite{*}.


%%% Put references here
\putbib
\end{bibunit}



\printglossary[type=\acronymtype]

\end{document}
